\documentclass[]{article}
\usepackage{lmodern}
\usepackage{amssymb,amsmath}
\usepackage{ifxetex,ifluatex}
\usepackage{fixltx2e} % provides \textsubscript
\ifnum 0\ifxetex 1\fi\ifluatex 1\fi=0 % if pdftex
  \usepackage[T1]{fontenc}
  \usepackage[utf8]{inputenc}
\else % if luatex or xelatex
  \ifxetex
    \usepackage{mathspec}
    \usepackage{xltxtra,xunicode}
  \else
    \usepackage{fontspec}
  \fi
  \defaultfontfeatures{Mapping=tex-text,Scale=MatchLowercase}
  \newcommand{\euro}{€}
\fi
% use upquote if available, for straight quotes in verbatim environments
\IfFileExists{upquote.sty}{\usepackage{upquote}}{}
% use microtype if available
\IfFileExists{microtype.sty}{\usepackage{microtype}}{}
\usepackage[margin=1in]{geometry}
\usepackage{color}
\usepackage{fancyvrb}
\newcommand{\VerbBar}{|}
\newcommand{\VERB}{\Verb[commandchars=\\\{\}]}
\DefineVerbatimEnvironment{Highlighting}{Verbatim}{commandchars=\\\{\}}
% Add ',fontsize=\small' for more characters per line
\newenvironment{Shaded}{}{}
\newcommand{\AlertTok}[1]{\textcolor[rgb]{1.00,0.00,0.00}{\textbf{#1}}}
\newcommand{\AnnotationTok}[1]{\textcolor[rgb]{0.38,0.63,0.69}{\textbf{\textit{#1}}}}
\newcommand{\AttributeTok}[1]{\textcolor[rgb]{0.49,0.56,0.16}{#1}}
\newcommand{\BaseNTok}[1]{\textcolor[rgb]{0.25,0.63,0.44}{#1}}
\newcommand{\BuiltInTok}[1]{\textcolor[rgb]{0.00,0.50,0.00}{#1}}
\newcommand{\CharTok}[1]{\textcolor[rgb]{0.25,0.44,0.63}{#1}}
\newcommand{\CommentTok}[1]{\textcolor[rgb]{0.38,0.63,0.69}{\textit{#1}}}
\newcommand{\CommentVarTok}[1]{\textcolor[rgb]{0.38,0.63,0.69}{\textbf{\textit{#1}}}}
\newcommand{\ConstantTok}[1]{\textcolor[rgb]{0.53,0.00,0.00}{#1}}
\newcommand{\ControlFlowTok}[1]{\textcolor[rgb]{0.00,0.44,0.13}{\textbf{#1}}}
\newcommand{\DataTypeTok}[1]{\textcolor[rgb]{0.56,0.13,0.00}{#1}}
\newcommand{\DecValTok}[1]{\textcolor[rgb]{0.25,0.63,0.44}{#1}}
\newcommand{\DocumentationTok}[1]{\textcolor[rgb]{0.73,0.13,0.13}{\textit{#1}}}
\newcommand{\ErrorTok}[1]{\textcolor[rgb]{1.00,0.00,0.00}{\textbf{#1}}}
\newcommand{\ExtensionTok}[1]{#1}
\newcommand{\FloatTok}[1]{\textcolor[rgb]{0.25,0.63,0.44}{#1}}
\newcommand{\FunctionTok}[1]{\textcolor[rgb]{0.02,0.16,0.49}{#1}}
\newcommand{\ImportTok}[1]{\textcolor[rgb]{0.00,0.50,0.00}{\textbf{#1}}}
\newcommand{\InformationTok}[1]{\textcolor[rgb]{0.38,0.63,0.69}{\textbf{\textit{#1}}}}
\newcommand{\KeywordTok}[1]{\textcolor[rgb]{0.00,0.44,0.13}{\textbf{#1}}}
\newcommand{\NormalTok}[1]{#1}
\newcommand{\OperatorTok}[1]{\textcolor[rgb]{0.40,0.40,0.40}{#1}}
\newcommand{\OtherTok}[1]{\textcolor[rgb]{0.00,0.44,0.13}{#1}}
\newcommand{\PreprocessorTok}[1]{\textcolor[rgb]{0.74,0.48,0.00}{#1}}
\newcommand{\RegionMarkerTok}[1]{#1}
\newcommand{\SpecialCharTok}[1]{\textcolor[rgb]{0.25,0.44,0.63}{#1}}
\newcommand{\SpecialStringTok}[1]{\textcolor[rgb]{0.73,0.40,0.53}{#1}}
\newcommand{\StringTok}[1]{\textcolor[rgb]{0.25,0.44,0.63}{#1}}
\newcommand{\VariableTok}[1]{\textcolor[rgb]{0.10,0.09,0.49}{#1}}
\newcommand{\VerbatimStringTok}[1]{\textcolor[rgb]{0.25,0.44,0.63}{#1}}
\newcommand{\WarningTok}[1]{\textcolor[rgb]{0.38,0.63,0.69}{\textbf{\textit{#1}}}}
\usepackage{graphicx}
\makeatletter
\def\maxwidth{\ifdim\Gin@nat@width>\linewidth\linewidth\else\Gin@nat@width\fi}
\def\maxheight{\ifdim\Gin@nat@height>\textheight\textheight\else\Gin@nat@height\fi}
\makeatother
% Scale images if necessary, so that they will not overflow the page
% margins by default, and it is still possible to overwrite the defaults
% using explicit options in \includegraphics[width, height, ...]{}
\setkeys{Gin}{width=\maxwidth,height=\maxheight,keepaspectratio}
\ifxetex
  \usepackage[setpagesize=false, % page size defined by xetex
              unicode=false, % unicode breaks when used with xetex
              xetex]{hyperref}
\else
  \usepackage[unicode=true]{hyperref}
\fi
\hypersetup{breaklinks=true,
            bookmarks=true,
            pdfauthor={Justin Le},
            pdftitle={Extreme Haskell: Typed State Machines (Part 2)},
            colorlinks=true,
            citecolor=blue,
            urlcolor=blue,
            linkcolor=magenta,
            pdfborder={0 0 0}}
\urlstyle{same}  % don't use monospace font for urls
% Make links footnotes instead of hotlinks:
\renewcommand{\href}[2]{#2\footnote{\url{#1}}}
\setlength{\parindent}{0pt}
\setlength{\parskip}{6pt plus 2pt minus 1pt}
\setlength{\emergencystretch}{3em}  % prevent overfull lines
\setcounter{secnumdepth}{0}

\title{Extreme Haskell: Typed State Machines (Part 2)}
\author{Justin Le}

\begin{document}
\maketitle

\emph{Originally posted on
\textbf{\href{https://blog.jle.im/entry/extreme-haskell-typed-state-machines.html}{in
Code}}.}

\documentclass[]{}
\usepackage{lmodern}
\usepackage{amssymb,amsmath}
\usepackage{ifxetex,ifluatex}
\usepackage{fixltx2e} % provides \textsubscript
\ifnum 0\ifxetex 1\fi\ifluatex 1\fi=0 % if pdftex
  \usepackage[T1]{fontenc}
  \usepackage[utf8]{inputenc}
\else % if luatex or xelatex
  \ifxetex
    \usepackage{mathspec}
    \usepackage{xltxtra,xunicode}
  \else
    \usepackage{fontspec}
  \fi
  \defaultfontfeatures{Mapping=tex-text,Scale=MatchLowercase}
  \newcommand{\euro}{€}
\fi
% use upquote if available, for straight quotes in verbatim environments
\IfFileExists{upquote.sty}{\usepackage{upquote}}{}
% use microtype if available
\IfFileExists{microtype.sty}{\usepackage{microtype}}{}
\usepackage[margin=1in]{geometry}
\ifxetex
  \usepackage[setpagesize=false, % page size defined by xetex
              unicode=false, % unicode breaks when used with xetex
              xetex]{hyperref}
\else
  \usepackage[unicode=true]{hyperref}
\fi
\hypersetup{breaklinks=true,
            bookmarks=true,
            pdfauthor={},
            pdftitle={},
            colorlinks=true,
            citecolor=blue,
            urlcolor=blue,
            linkcolor=magenta,
            pdfborder={0 0 0}}
\urlstyle{same}  % don't use monospace font for urls
% Make links footnotes instead of hotlinks:
\renewcommand{\href}[2]{#2\footnote{\url{#1}}}
\setlength{\parindent}{0pt}
\setlength{\parskip}{6pt plus 2pt minus 1pt}
\setlength{\emergencystretch}{3em}  % prevent overfull lines
\setcounter{secnumdepth}{0}

\textbackslash begin\{document\}

Welcome back to our journey through the limits of polite Haskell! In
\href{/entry/extreme-haskell-typed-expression-edsls-1.html}{Part 1}, we built a
typed expression language whose terms track their free variables and result
types, and whose evaluator is total. If GHC lets you construct the
\texttt{Expr}, it will evaluate.

That is already useful, but I would not have written a thousand lines about it
if the payoff was just ``a lambda calculus that typechecks.'' The real payoff is
that the expression language is now a reusable typed component. You can embed it
inside larger typed structures, and the type indices thread through, carrying
their guarantees with them. In this post, we will put it inside a typed
state-machine description: a graph of states and triggered transitions where the
guards, the payloads, and the destinations are all verified at compile time. We
will then render those machines as Graphviz diagrams directly from the typed
structure.

This is where I personally started feeling the approach justify itself. You
write one machine definition, and you get a Graphviz diagram for free, a
JavaScript simulator for free, and the guarantee that they all agree with each
other for free. When I first got this working on a real project, the feeling was
less ``oh neat'' and more ``why would anyone do this any other way.''

Again, all of the code here is
\href{https://github.com/mstksg/inCode/tree/master/code-samples/typed-sm-lc/flake.nix}{available
online}, and you can load it all in ghci from the sample directory:

\begin{Shaded}
\begin{Highlighting}[]
\ExtensionTok{$}\NormalTok{ cd code{-}samples/typed{-}sm{-}lc}
\ExtensionTok{$}\NormalTok{ nix develop}
\ExtensionTok{$}\NormalTok{ ghci}
\ExtensionTok{ghci}\OperatorTok{\textgreater{}}\NormalTok{ :load MachineStage5.hs}
\end{Highlighting}
\end{Shaded}

\section{The State Machine}\label{the-state-machine}

Okay, now here is where all of this finally starts paying rent.

The expression language we made is not merely something that evaluates to a
value in Haskell. It is a typed description of a computation. And if it is a
description, we can stick it inside a larger typed description: a state machine.

\subsection{A Machine Made of Strings}\label{a-machine-made-of-strings}

For our running example, we will build a controller for a software-defined radio
(SDR) receiver. An SDR is a radio where the signal processing that traditionally
happens in hardware (tuning, filtering, demodulation) is instead done in
software. The receiver scans across a band of frequencies looking for signals of
interest, locks onto them when found, and handles the various failure modes
(signal drift, timeout, thermal limits) that come up in practice.

This is a natural fit for a state machine. The controller has a small number of
discrete modes (idle, scanning, locked onto a signal, cooling down), and the
transitions between them depend on real-time sensor readings (received signal
strength, noise floor, frequency offset). Each mode carries different data: a
scanning receiver tracks its current frequency and how long it has dwelt there;
a locked receiver additionally knows the signal strength at lock time. The
transitions have guards (``only lock if the signal is strong enough'') and must
construct the right payload for the destination state.

State machines like this show up everywhere in embedded systems, protocol
implementations, and game logic. They are also where I have personally seen the
most insidious runtime bugs: transitions that reference data from the wrong
state, destinations with incomplete payloads, guards that accidentally test the
wrong condition. The typed version we build here eliminates all three classes of
bug at compile time.

So: we have a few states (idle, scanning, locked, cooling down), some triggers
(start, peak detected, drift, timeout, reset), and some transitions between them
with guards and payload constructors.

An obvious first pass in Haskell is something like:

\begin{Shaded}
\begin{Highlighting}[]
\CommentTok{{-}{-} source: https://github.com/mstksg/inCode/tree/master/code{-}samples/typed{-}sm{-}lc/MachineUntyped.hs\#L8{-}L16}

\KeywordTok{data} \DataTypeTok{UntypedTransition} \OtherTok{=} \DataTypeTok{UntypedTransition}
\NormalTok{  \{}\OtherTok{ utFrom    ::} \DataTypeTok{String}
\NormalTok{  ,}\OtherTok{ utTrigger ::} \DataTypeTok{String}
\NormalTok{  ,}\OtherTok{ utTo      ::} \DataTypeTok{String}
\NormalTok{  ,}\OtherTok{ utGuard   ::} \DataTypeTok{Map} \DataTypeTok{String} \DataTypeTok{Int} \OtherTok{{-}\textgreater{}} \DataTypeTok{Bool}
\NormalTok{  ,}\OtherTok{ utBuild   ::} \DataTypeTok{Map} \DataTypeTok{String} \DataTypeTok{Int} \OtherTok{{-}\textgreater{}} \DataTypeTok{Map} \DataTypeTok{String} \DataTypeTok{Int}
\NormalTok{  \}}

\KeywordTok{type} \DataTypeTok{UntypedMachine} \OtherTok{=}\NormalTok{ [}\DataTypeTok{UntypedTransition}\NormalTok{]}
\end{Highlighting}
\end{Shaded}

States are strings. Triggers are strings. The payload is a bag of string-keyed
integers. Guards are opaque Haskell functions over that bag. This is pleasant to
write! You can throw a machine together in five minutes, and there is zero
ceremony around declaring your state space or its structure.

Here is a complete machine in this style:

\begin{Shaded}
\begin{Highlighting}[]
\CommentTok{{-}{-} source: https://github.com/mstksg/inCode/tree/master/code{-}samples/typed{-}sm{-}lc/MachineUntyped.hs\#L37{-}L61}

\OtherTok{scanMachine ::} \DataTypeTok{UntypedMachine}
\NormalTok{scanMachine }\OtherTok{=}
\NormalTok{  [ }\DataTypeTok{UntypedTransition} \StringTok{"Idle"} \StringTok{"Start"} \StringTok{"Scanning"}
\NormalTok{      (}\FunctionTok{const} \DataTypeTok{True}\NormalTok{)}
\NormalTok{      (\textbackslash{}\_ }\OtherTok{{-}\textgreater{}}\NormalTok{ M.fromList [(}\StringTok{"freq"}\NormalTok{, }\DecValTok{2400}\NormalTok{), (}\StringTok{"dwell"}\NormalTok{, }\DecValTok{0}\NormalTok{)])}
\NormalTok{  , }\DataTypeTok{UntypedTransition} \StringTok{"Scanning"} \StringTok{"Peak"} \StringTok{"Locked"}
\NormalTok{      (\textbackslash{}env }\OtherTok{{-}\textgreater{}}\NormalTok{ (env }\OperatorTok{M.!} \StringTok{"rssi"}\NormalTok{) }\OperatorTok{\textgreater{}=}\NormalTok{ (env }\OperatorTok{M.!} \StringTok{"noise"}\NormalTok{) }\OperatorTok{+} \DecValTok{4}\NormalTok{)}
\NormalTok{      (\textbackslash{}env }\OtherTok{{-}\textgreater{}}\NormalTok{ M.fromList}
\NormalTok{        [ (}\StringTok{"freq"}\NormalTok{, (env }\OperatorTok{M.!} \StringTok{"freq"}\NormalTok{) }\OperatorTok{+}\NormalTok{ (env }\OperatorTok{M.!} \StringTok{"peakOffset"}\NormalTok{))}
\NormalTok{        , (}\StringTok{"lockRssi"}\NormalTok{, env }\OperatorTok{M.!} \StringTok{"rssi"}\NormalTok{)}
\NormalTok{        , (}\StringTok{"dwell"}\NormalTok{, (env }\OperatorTok{M.!} \StringTok{"dwell"}\NormalTok{) }\OperatorTok{+} \DecValTok{1}\NormalTok{)}
\NormalTok{        ])}
\NormalTok{  , }\DataTypeTok{UntypedTransition} \StringTok{"Scanning"} \StringTok{"Timeout"} \StringTok{"CoolingDown"}
\NormalTok{      (\textbackslash{}env }\OtherTok{{-}\textgreater{}}\NormalTok{ (env }\OperatorTok{M.!} \StringTok{"dwell"}\NormalTok{) }\OperatorTok{\textgreater{}=} \DecValTok{3}\NormalTok{)}
\NormalTok{      (\textbackslash{}env }\OtherTok{{-}\textgreater{}}\NormalTok{ M.singleton }\StringTok{"lastFreq"}\NormalTok{ (env }\OperatorTok{M.!} \StringTok{"freq"}\NormalTok{))}
\NormalTok{  , }\DataTypeTok{UntypedTransition} \StringTok{"Locked"} \StringTok{"Drift"} \StringTok{"Scanning"}
\NormalTok{      (}\FunctionTok{const} \DataTypeTok{True}\NormalTok{)}
\NormalTok{      (\textbackslash{}env }\OtherTok{{-}\textgreater{}}\NormalTok{ M.fromList}
\NormalTok{        [ (}\StringTok{"freq"}\NormalTok{, (env }\OperatorTok{M.!} \StringTok{"freq"}\NormalTok{) }\OperatorTok{+}\NormalTok{ (env }\OperatorTok{M.!} \StringTok{"peakOffset"}\NormalTok{))}
\NormalTok{        , (}\StringTok{"dwell"}\NormalTok{, env }\OperatorTok{M.!} \StringTok{"dwell"}\NormalTok{)}
\NormalTok{        ])}
\NormalTok{  , }\DataTypeTok{UntypedTransition} \StringTok{"CoolingDown"} \StringTok{"Reset"} \StringTok{"Idle"}
\NormalTok{      (}\FunctionTok{const} \DataTypeTok{True}\NormalTok{)}
\NormalTok{      (}\FunctionTok{const}\NormalTok{ M.empty)}
\NormalTok{  ]}
\end{Highlighting}
\end{Shaded}

Looks reasonable. But now consider:

\begin{Shaded}
\begin{Highlighting}[]
\CommentTok{{-}{-} source: https://github.com/mstksg/inCode/tree/master/code{-}samples/typed{-}sm{-}lc/MachineUntyped.hs\#L18{-}L25}

\OtherTok{badTransition1 ::} \DataTypeTok{UntypedTransition}
\NormalTok{badTransition1 }\OtherTok{=} \DataTypeTok{UntypedTransition}
\NormalTok{  \{ utFrom    }\OtherTok{=} \StringTok{"Scanning"}
\NormalTok{  , utTrigger }\OtherTok{=} \StringTok{"Peak"}
\NormalTok{  , utTo      }\OtherTok{=} \StringTok{"Lockde"}          \CommentTok{{-}{-} typo: "Lockde" instead of "Locked"}
\NormalTok{  , utGuard   }\OtherTok{=}\NormalTok{ \textbackslash{}env }\OtherTok{{-}\textgreater{}}\NormalTok{ (env }\OperatorTok{M.!} \StringTok{"rssi"}\NormalTok{) }\OperatorTok{\textgreater{}} \DecValTok{4}
\NormalTok{  , utBuild   }\OtherTok{=}\NormalTok{ \textbackslash{}env }\OtherTok{{-}\textgreater{}}\NormalTok{ M.singleton }\StringTok{"lockRssi"}\NormalTok{ (env }\OperatorTok{M.!} \StringTok{"rssi"}\NormalTok{)}
\NormalTok{  \}}
\end{Highlighting}
\end{Shaded}

This transitions to \texttt{"Lockde"} instead of \texttt{"Locked"}. GHC sees two
perfectly valid \texttt{String} values and has nothing to say about it. The typo
is only observable at runtime, when the machine steps into a state that has no
outgoing edges.

\begin{Shaded}
\begin{Highlighting}[]
\CommentTok{{-}{-} source: https://github.com/mstksg/inCode/tree/master/code{-}samples/typed{-}sm{-}lc/MachineUntyped.hs\#L27{-}L35}

\OtherTok{badTransition2 ::} \DataTypeTok{UntypedTransition}
\NormalTok{badTransition2 }\OtherTok{=} \DataTypeTok{UntypedTransition}
\NormalTok{  \{ utFrom    }\OtherTok{=} \StringTok{"Scanning"}
\NormalTok{  , utTrigger }\OtherTok{=} \StringTok{"Peak"}
\NormalTok{  , utTo      }\OtherTok{=} \StringTok{"Locked"}
\NormalTok{  , utGuard   }\OtherTok{=}\NormalTok{ \textbackslash{}env }\OtherTok{{-}\textgreater{}}\NormalTok{ (env }\OperatorTok{M.!} \StringTok{"lockRssi"}\NormalTok{) }\OperatorTok{\textgreater{}} \DecValTok{4}  \CommentTok{{-}{-} lockRssi doesn\textquotesingle{}t exist while Scanning!}
\NormalTok{  , utBuild   }\OtherTok{=}\NormalTok{ \textbackslash{}env }\OtherTok{{-}\textgreater{}}\NormalTok{ M.singleton }\StringTok{"freq"}\NormalTok{ (env }\OperatorTok{M.!} \StringTok{"rssi"}\NormalTok{)}
                        \CommentTok{{-}{-} missing "lockRssi" and "dwell" that Locked needs}
\NormalTok{  \}}
\end{Highlighting}
\end{Shaded}

This one reads \texttt{"lockRssi"} from the environment while in the
\texttt{"Scanning"} state. But \texttt{lockRssi} only exists after locking on to
a signal. It also builds only \texttt{"freq"} for the destination payload, but
\texttt{"Locked"} needs \texttt{"freq"}, \texttt{"lockRssi"}, and
\texttt{"dwell"}. Both compile. Both pass any test that doesn't exercise exactly
these edges. Both explode at runtime with a key lookup failure.

This is the state-machine version of \texttt{1\ \&\&\ 2} from Part 1. The system
cannot distinguish a well-formed transition from garbage because the types carry
no information about which states exist, what variables they carry, or what
shape the destination payload should have.

And unlike \texttt{1\ \&\&\ 2}, which you would catch immediately in testing, a
bad state machine transition can hide for months. In my experience, state
machine bugs are the worst kind: they only manifest under specific sequences of
events, they are hard to reproduce, and by the time you see the failure in
production the causal chain is three state transitions deep. I have personally
debugged enough of these that I refuse to write another untyped state machine if
I can help it.

We are going to fix it the same way as Part 1: by making the illegal versions
un-representable.

\subsection{Closing the World of States}\label{closing-the-world-of-states}

The first repair is the same one we applied to \texttt{Ty} in Part 1: instead of
letting any string be a state or trigger, we define a closed kind.

\begin{Shaded}
\begin{Highlighting}[]
\CommentTok{{-}{-} source: https://github.com/mstksg/inCode/tree/master/code{-}samples/typed{-}sm{-}lc/MachineStage5.hs\#L22{-}L33}

\KeywordTok{type} \KeywordTok{data} \DataTypeTok{ScanState}
  \OtherTok{=} \DataTypeTok{Idle}
  \OperatorTok{|} \DataTypeTok{Scanning}
  \OperatorTok{|} \DataTypeTok{Locked}
  \OperatorTok{|} \DataTypeTok{CoolingDown}

\KeywordTok{type} \KeywordTok{data} \DataTypeTok{Trigger}
  \OtherTok{=} \DataTypeTok{Start}
  \OperatorTok{|} \DataTypeTok{Peak}
  \OperatorTok{|} \DataTypeTok{Drift}
  \OperatorTok{|} \DataTypeTok{Timeout}
  \OperatorTok{|} \DataTypeTok{Reset}
\end{Highlighting}
\end{Shaded}

\texttt{ScanState} is a kind with exactly four types: \texttt{Idle},
\texttt{Scanning}, \texttt{Locked}, \texttt{CoolingDown}. \texttt{Trigger} is a
kind with five: \texttt{Start}, \texttt{Peak}, \texttt{Drift}, \texttt{Timeout},
\texttt{Reset}. These are declared with \texttt{type\ data}, the same mechanism
we used for \texttt{Ty} in Part 1. Remember: \texttt{type\ data} creates types
that exist only at the type level. There is no value-level \texttt{Idle}
constructor floating around. \texttt{Idle} is a \emph{type} of kind
\texttt{ScanState}, and that is all it is.

The singletons \texttt{SState} and \texttt{STrigger} are the runtime handles we
keep around so we can render state and trigger names later (for Graphviz labels,
JavaScript codegen, etc.):

\begin{Shaded}
\begin{Highlighting}[]
\CommentTok{{-}{-} source: https://github.com/mstksg/inCode/tree/master/code{-}samples/typed{-}sm{-}lc/MachineStage5.hs\#L49{-}L60}

\KeywordTok{data} \DataTypeTok{SState}\OtherTok{ ::} \DataTypeTok{ScanState} \OtherTok{{-}\textgreater{}} \DataTypeTok{Type} \KeywordTok{where}
  \DataTypeTok{SIdle}\OtherTok{ ::} \DataTypeTok{SState} \DataTypeTok{Idle}
  \DataTypeTok{SScanning}\OtherTok{ ::} \DataTypeTok{SState} \DataTypeTok{Scanning}
  \DataTypeTok{SLocked}\OtherTok{ ::} \DataTypeTok{SState} \DataTypeTok{Locked}
  \DataTypeTok{SCoolingDown}\OtherTok{ ::} \DataTypeTok{SState} \DataTypeTok{CoolingDown}

\KeywordTok{data} \DataTypeTok{STrigger}\OtherTok{ ::} \DataTypeTok{Trigger} \OtherTok{{-}\textgreater{}} \DataTypeTok{Type} \KeywordTok{where}
  \DataTypeTok{SStart}\OtherTok{ ::} \DataTypeTok{STrigger} \DataTypeTok{Start}
  \DataTypeTok{SPeak}\OtherTok{ ::} \DataTypeTok{STrigger} \DataTypeTok{Peak}
  \DataTypeTok{SDrift}\OtherTok{ ::} \DataTypeTok{STrigger} \DataTypeTok{Drift}
  \DataTypeTok{STimeout}\OtherTok{ ::} \DataTypeTok{STrigger} \DataTypeTok{Timeout}
  \DataTypeTok{SReset}\OtherTok{ ::} \DataTypeTok{STrigger} \DataTypeTok{Reset}
\end{Highlighting}
\end{Shaded}

This is exactly the same pattern as \texttt{STy} from Part 1, just applied to a
different domain. \texttt{SState} has one constructor per inhabitant of
\texttt{ScanState}, each one indexed by the corresponding type. Pattern matching
on \texttt{SIdle} recovers the knowledge \texttt{s\ \textasciitilde{}\ Idle}.

We also define name functions that extract the string label from each singleton,
which we will use later for rendering:

\begin{Shaded}
\begin{Highlighting}[]
\CommentTok{{-}{-} source: https://github.com/mstksg/inCode/tree/master/code{-}samples/typed{-}sm{-}lc/MachineStage5.hs\#L62{-}L75}

\OtherTok{stateName ::} \DataTypeTok{SState}\NormalTok{ s }\OtherTok{{-}\textgreater{}} \DataTypeTok{String}
\NormalTok{stateName }\OtherTok{=}\NormalTok{ \textbackslash{}}\KeywordTok{case}
  \DataTypeTok{SIdle} \OtherTok{{-}\textgreater{}} \StringTok{"Idle"}
  \DataTypeTok{SScanning} \OtherTok{{-}\textgreater{}} \StringTok{"Scanning"}
  \DataTypeTok{SLocked} \OtherTok{{-}\textgreater{}} \StringTok{"Locked"}
  \DataTypeTok{SCoolingDown} \OtherTok{{-}\textgreater{}} \StringTok{"CoolingDown"}

\OtherTok{triggerName ::} \DataTypeTok{STrigger}\NormalTok{ t }\OtherTok{{-}\textgreater{}} \DataTypeTok{String}
\NormalTok{triggerName }\OtherTok{=}\NormalTok{ \textbackslash{}}\KeywordTok{case}
  \DataTypeTok{SStart} \OtherTok{{-}\textgreater{}} \StringTok{"Start"}
  \DataTypeTok{SPeak} \OtherTok{{-}\textgreater{}} \StringTok{"Peak"}
  \DataTypeTok{SDrift} \OtherTok{{-}\textgreater{}} \StringTok{"Drift"}
  \DataTypeTok{STimeout} \OtherTok{{-}\textgreater{}} \StringTok{"Timeout"}
  \DataTypeTok{SReset} \OtherTok{{-}\textgreater{}} \StringTok{"Reset"}
\end{Highlighting}
\end{Shaded}

A transition from \texttt{"Lockde"} is now a type error: \texttt{Lockde} simply
does not exist in the \texttt{ScanState} kind. If you mistype it, GHC will tell
you immediately that there is no such type, rather than silently building a
value that refers to a nonexistent state.

\subsection{State Payloads Are Not All the
Same}\label{state-payloads-are-not-all-the-same}

The untyped machine stored everything in a \texttt{Map\ String\ Int}. But not
every state carries the same data. An idle controller has no payload at all. A
scanning controller tracks its current frequency and dwell count. A locked
controller additionally knows the signal strength at which it locked. And a
cooling-down controller only remembers the last frequency it was on.

In Part 1, we used type-level lists to track the variables in scope:
\texttt{Expr\ {[}"x"\ :::\ TInt,\ "y"\ :::\ TBool{]}\ TInt} meant ``an integer
expression that can reference \texttt{x} (an int) and \texttt{y} (a bool).'' The
variable environment was \emph{explicit} in the type. We want the same thing
here, but now the variable environment depends on \emph{which state we are in}.
We need a way to compute a type-level list from a type-level state.

This is what a \emph{type family} does. A type family is a function at the type
level: given a type as input, it produces a type as output. We have not needed
one until now because in Part 1 the scope was always provided directly. Here,
the scope is \emph{derived} from the state:

\begin{Shaded}
\begin{Highlighting}[]
\CommentTok{{-}{-} source: https://github.com/mstksg/inCode/tree/master/code{-}samples/typed{-}sm{-}lc/MachineStage5.hs\#L35{-}L47}

\KeywordTok{type} \KeywordTok{family} \DataTypeTok{NodeVars}\NormalTok{ (}\OtherTok{s ::} \DataTypeTok{ScanState}\NormalTok{) }\KeywordTok{where}
  \DataTypeTok{NodeVars} \DataTypeTok{Idle} \OtherTok{=}\NormalTok{ \textquotesingle{}[]}
  \DataTypeTok{NodeVars} \DataTypeTok{Scanning} \OtherTok{=}\NormalTok{ \textquotesingle{}[}\StringTok{"freq"} \OperatorTok{:::} \DataTypeTok{TInt}\NormalTok{, }\StringTok{"dwell"} \OperatorTok{:::} \DataTypeTok{TInt}\NormalTok{]}
  \DataTypeTok{NodeVars} \DataTypeTok{Locked} \OtherTok{=}\NormalTok{ \textquotesingle{}[}\StringTok{"freq"} \OperatorTok{:::} \DataTypeTok{TInt}\NormalTok{, }\StringTok{"lockRssi"} \OperatorTok{:::} \DataTypeTok{TInt}\NormalTok{, }\StringTok{"dwell"} \OperatorTok{:::} \DataTypeTok{TInt}\NormalTok{]}
  \DataTypeTok{NodeVars} \DataTypeTok{CoolingDown} \OtherTok{=}\NormalTok{ \textquotesingle{}[}\StringTok{"lastFreq"} \OperatorTok{:::} \DataTypeTok{TInt}\NormalTok{]}

\KeywordTok{type} \DataTypeTok{Env} \OtherTok{=}
\NormalTok{  \textquotesingle{}[ }\StringTok{"rssi"} \OperatorTok{:::} \DataTypeTok{TInt}\NormalTok{,}
     \StringTok{"noise"} \OperatorTok{:::} \DataTypeTok{TInt}\NormalTok{,}
     \StringTok{"peakOffset"} \OperatorTok{:::} \DataTypeTok{TInt}
\NormalTok{   ]}

\KeywordTok{type} \DataTypeTok{Scope}\NormalTok{ s }\OtherTok{=} \DataTypeTok{Env} \OperatorTok{++} \DataTypeTok{NodeVars}\NormalTok{ s}
\end{Highlighting}
\end{Shaded}

\texttt{NodeVars} is a closed type family: you give it a \texttt{ScanState} and
it returns the corresponding variable list. \texttt{NodeVars\ Idle} reduces to
\texttt{{[}{]}}. \texttt{NodeVars\ Scanning} reduces to
\texttt{{[}"freq"\ :::\ TInt,\ "dwell"\ :::\ TInt{]}}. \texttt{NodeVars\ Locked}
reduces to
\texttt{{[}"freq"\ :::\ TInt,\ "lockRssi"\ :::\ TInt,\ "dwell"\ :::\ TInt{]}}.
\texttt{NodeVars\ CoolingDown} reduces to \texttt{{[}"lastFreq"\ :::\ TInt{]}}.
These reductions happen at compile time: GHC evaluates the type family when it
needs to check whether two types are equal.

The ambient \texttt{Env} is the external observations (rssi, noise floor, peak
offset) that are always available regardless of the current state. These are
things the machine reads from its environment on every tick. \texttt{Env} is a
plain type synonym, not a type family, because it does not depend on anything.

And \texttt{Scope\ s\ =\ Env\ ++\ NodeVars\ s} is the \emph{full} variable
environment available to guards and builders while in state \texttt{s}. This is
the direct analog of the \texttt{vs} parameter in \texttt{Expr\ vs\ t} from Part
1. A guard expression in state \texttt{Scanning} has type
\texttt{Expr\ (Scope\ Scanning)\ TBool}, which means it can reference
\texttt{rssi}, \texttt{noise}, \texttt{peakOffset}, \texttt{freq}, and
\texttt{dwell}. It \emph{cannot} reference \texttt{lockRssi}, because that
variable only exists in \texttt{Scope\ Locked}.

We need a type-level append to concatenate the environment and state-local
variables:

\begin{Shaded}
\begin{Highlighting}[]
\CommentTok{{-}{-} source: https://github.com/mstksg/inCode/tree/master/code{-}samples/typed{-}sm{-}lc/MachineStage5.hs\#L18{-}L20}

\KeywordTok{type} \KeywordTok{family}\NormalTok{ xs }\OperatorTok{++}\NormalTok{ ys }\KeywordTok{where}
\NormalTok{  \textquotesingle{}[] }\OperatorTok{++}\NormalTok{ ys }\OtherTok{=}\NormalTok{ ys}
\NormalTok{  (x \textquotesingle{}}\OperatorTok{:}\NormalTok{ xs) }\OperatorTok{++}\NormalTok{ ys }\OtherTok{=}\NormalTok{ x \textquotesingle{}}\OperatorTok{:}\NormalTok{ (xs }\OperatorTok{++}\NormalTok{ ys)}
\end{Highlighting}
\end{Shaded}

Another closed type family. It peels elements off the left list and conses them
onto the right, exactly like \texttt{(++)} at the value level but operating on
promoted lists. Let us trace through the reduction manually for
\texttt{Scope\ Scanning}:

\begin{verbatim}
Scope Scanning
  = Env ++ NodeVars Scanning
  = '["rssi" ::: TInt, "noise" ::: TInt, "peakOffset" ::: TInt]
      ++ '["freq" ::: TInt, "dwell" ::: TInt]
  = '[ "rssi" ::: TInt, "noise" ::: TInt, "peakOffset" ::: TInt
     , "freq" ::: TInt, "dwell" ::: TInt ]
\end{verbatim}

The result is a flat type-level list of five named, typed variables. This is the
scope that guards and builders operate over when the machine is in the
\texttt{Scanning} state. Compare with \texttt{Scope\ Idle}:

\begin{verbatim}
Scope Idle
  = Env ++ NodeVars Idle
  = Env ++ '[]
  = Env
  = '["rssi" ::: TInt, "noise" ::: TInt, "peakOffset" ::: TInt]
\end{verbatim}

Only three variables. A guard in the \texttt{Idle} state can reference the
environment but nothing else, because there is no state-local payload. And
\texttt{Scope\ Locked}:

\begin{verbatim}
Scope Locked
  = Env ++ NodeVars Locked
  = '[ "rssi" ::: TInt, "noise" ::: TInt, "peakOffset" ::: TInt
     , "freq" ::: TInt, "lockRssi" ::: TInt, "dwell" ::: TInt ]
\end{verbatim}

Six variables. The key insight: these are all just type-level lists with the
same \texttt{(Symbol,\ Ty)} element type. They plug directly into the
\texttt{Expr\ vs\ t} and \texttt{Rec\ f\ vs} machinery from Part 1 with zero
additional infrastructure. The expression language already knows how to work
with any type-level list of named typed variables. We just need to compute the
right list for each state.

If you try to use a \texttt{Scope\ Locked} expression as a guard while in
\texttt{Scanning}, GHC will reject it with a type mismatch on the variable
environment. The ``reading a variable from the wrong state'' class of bug is
gone.

\subsection{A Guard Is Just an Expression}\label{a-guard-is-just-an-expression}

Let us build up a concrete guard step by step. While scanning, we want to lock
onto a signal if it is strong enough: specifically, if
\texttt{rssi\ \textgreater{}=\ noise\ +\ threshold}. In our expression language,
we need to reference variables from \texttt{Scope\ Scanning}.

Recall from Part 1 that variable references are de Bruijn-style indices into the
scope. An \texttt{EVar} takes an \texttt{Index\ vs\ (name\ :::\ t)}, which
counts how many steps from the head of the list we need to skip. The scope for
\texttt{Scanning} is:

\begin{verbatim}
Scope Scanning = Env ++ NodeVars Scanning
               = '[ "rssi"       ::: TInt     -- position 0 (IZ)
                  , "noise"      ::: TInt     -- position 1 (IS IZ)
                  , "peakOffset" ::: TInt     -- position 2 (IS (IS IZ))
                  , "freq"       ::: TInt     -- position 3 (IS (IS (IS IZ)))
                  , "dwell"      ::: TInt     -- position 4 (IS (IS (IS (IS IZ))))
                  ]
\end{verbatim}

So to reference \texttt{rssi} while scanning, we write \texttt{EVar\ IZ}. To
reference \texttt{noise}, we write \texttt{EVar\ (IS\ IZ)}. And so on. We define
named aliases so we do not have to count indices by hand every time:

\begin{Shaded}
\begin{Highlighting}[]
\CommentTok{{-}{-} source: https://github.com/mstksg/inCode/tree/master/code{-}samples/typed{-}sm{-}lc/MachineStage5.hs\#L149{-}L156}

\OtherTok{rssiScanning ::} \DataTypeTok{Expr}\NormalTok{ (}\DataTypeTok{Scope} \DataTypeTok{Scanning}\NormalTok{) }\DataTypeTok{TInt}
\NormalTok{rssiScanning }\OtherTok{=} \DataTypeTok{EVar} \DataTypeTok{IZ}

\OtherTok{noiseScanning ::} \DataTypeTok{Expr}\NormalTok{ (}\DataTypeTok{Scope} \DataTypeTok{Scanning}\NormalTok{) }\DataTypeTok{TInt}
\NormalTok{noiseScanning }\OtherTok{=} \DataTypeTok{EVar}\NormalTok{ (}\DataTypeTok{IS} \DataTypeTok{IZ}\NormalTok{)}

\OtherTok{peakOffsetScanning ::} \DataTypeTok{Expr}\NormalTok{ (}\DataTypeTok{Scope} \DataTypeTok{Scanning}\NormalTok{) }\DataTypeTok{TInt}
\NormalTok{peakOffsetScanning }\OtherTok{=} \DataTypeTok{EVar}\NormalTok{ (}\DataTypeTok{IS}\NormalTok{ (}\DataTypeTok{IS} \DataTypeTok{IZ}\NormalTok{))}
\end{Highlighting}
\end{Shaded}

Each one is an \texttt{EVar} with a specific \texttt{Index} that picks out the
correct position in the scope. \texttt{rssiScanning} is \texttt{EVar\ IZ}
because \texttt{rssi} is the first variable in \texttt{Scope\ Scanning} (which
is \texttt{Env\ ++\ NodeVars\ Scanning}). \texttt{noiseScanning} is
\texttt{EVar\ (IS\ IZ)} because it's one step past the head.

Similarly for the \texttt{Scanning}-specific locals:

\begin{Shaded}
\begin{Highlighting}[]
\CommentTok{{-}{-} source: https://github.com/mstksg/inCode/tree/master/code{-}samples/typed{-}sm{-}lc/MachineStage5.hs\#L167{-}L171}

\OtherTok{scanFreq ::} \DataTypeTok{Expr}\NormalTok{ (}\DataTypeTok{Scope} \DataTypeTok{Scanning}\NormalTok{) }\DataTypeTok{TInt}
\NormalTok{scanFreq }\OtherTok{=} \DataTypeTok{EVar}\NormalTok{ (}\DataTypeTok{IS}\NormalTok{ (}\DataTypeTok{IS}\NormalTok{ (}\DataTypeTok{IS} \DataTypeTok{IZ}\NormalTok{)))}

\OtherTok{scanDwell ::} \DataTypeTok{Expr}\NormalTok{ (}\DataTypeTok{Scope} \DataTypeTok{Scanning}\NormalTok{) }\DataTypeTok{TInt}
\NormalTok{scanDwell }\OtherTok{=} \DataTypeTok{EVar}\NormalTok{ (}\DataTypeTok{IS}\NormalTok{ (}\DataTypeTok{IS}\NormalTok{ (}\DataTypeTok{IS}\NormalTok{ (}\DataTypeTok{IS} \DataTypeTok{IZ}\NormalTok{))))}
\end{Highlighting}
\end{Shaded}

\texttt{scanFreq} is at position 3 (\texttt{IS\ (IS\ (IS\ IZ))}): three
\texttt{IS} wrappers skip past the three \texttt{Env} entries, then \texttt{IZ}
lands on \texttt{"freq"}. \texttt{scanDwell} is at position 4.

Now we can write the guard itself:

\begin{Shaded}
\begin{Highlighting}[]
\CommentTok{{-}{-} source: https://github.com/mstksg/inCode/tree/master/code{-}samples/typed{-}sm{-}lc/MachineStage5.hs\#L197{-}L198}

\OtherTok{strongEnoughScanning ::} \DataTypeTok{Expr}\NormalTok{ (}\DataTypeTok{Scope} \DataTypeTok{Scanning}\NormalTok{) }\DataTypeTok{TBool}
\NormalTok{strongEnoughScanning }\OtherTok{=} \DataTypeTok{EOp} \DataTypeTok{OLte}\NormalTok{ (}\DataTypeTok{EOp} \DataTypeTok{OPlus}\NormalTok{ noiseScanning threshold) rssiScanning}
\end{Highlighting}
\end{Shaded}

This is \texttt{Expr\ (Scope\ Scanning)\ TBool}: a boolean expression whose free
variables are exactly those available while scanning. It says
\texttt{(noise\ +\ 4)\ \textless{}=\ rssi}. The \texttt{threshold} and other
constants are polymorphic in scope:

\begin{Shaded}
\begin{Highlighting}[]
\CommentTok{{-}{-} source: https://github.com/mstksg/inCode/tree/master/code{-}samples/typed{-}sm{-}lc/MachineStage5.hs\#L182{-}L186}

\OtherTok{threshold ::} \DataTypeTok{Expr}\NormalTok{ scope }\DataTypeTok{TInt}
\NormalTok{threshold }\OtherTok{=} \DataTypeTok{EPrim}\NormalTok{ (}\DataTypeTok{PInt} \DecValTok{4}\NormalTok{)}

\OtherTok{maxDwell ::} \DataTypeTok{Expr}\NormalTok{ scope }\DataTypeTok{TInt}
\NormalTok{maxDwell }\OtherTok{=} \DataTypeTok{EPrim}\NormalTok{ (}\DataTypeTok{PInt} \DecValTok{3}\NormalTok{)}
\end{Highlighting}
\end{Shaded}

Because a literal like \texttt{EPrim\ (PInt\ 4)} does not reference any
variables, it can live in any scope. So
\texttt{threshold\ ::\ Expr\ scope\ TInt} is parametrically polymorphic in the
scope, and can be used in any state's guard.

Because the guard must be \texttt{Expr\ (Scope\ s)\ TBool}, you get two
guarantees for free:

\begin{enumerate}
\def\labelenumi{\arabic{enumi}.}
\tightlist
\item
  It can only reference variables that exist in state \texttt{s}.
\item
  It must actually produce a boolean. You cannot accidentally use an integer
  expression as a transition predicate.
\end{enumerate}

Similarly for the \texttt{Locked} state, we have a separate set of variable
references into \texttt{Scope\ Locked}:

\begin{Shaded}
\begin{Highlighting}[]
\CommentTok{{-}{-} source: https://github.com/mstksg/inCode/tree/master/code{-}samples/typed{-}sm{-}lc/MachineStage5.hs\#L158{-}L165}

\OtherTok{rssiLocked ::} \DataTypeTok{Expr}\NormalTok{ (}\DataTypeTok{Scope} \DataTypeTok{Locked}\NormalTok{) }\DataTypeTok{TInt}
\NormalTok{rssiLocked }\OtherTok{=} \DataTypeTok{EVar} \DataTypeTok{IZ}

\OtherTok{noiseLocked ::} \DataTypeTok{Expr}\NormalTok{ (}\DataTypeTok{Scope} \DataTypeTok{Locked}\NormalTok{) }\DataTypeTok{TInt}
\NormalTok{noiseLocked }\OtherTok{=} \DataTypeTok{EVar}\NormalTok{ (}\DataTypeTok{IS} \DataTypeTok{IZ}\NormalTok{)}

\OtherTok{peakOffsetLocked ::} \DataTypeTok{Expr}\NormalTok{ (}\DataTypeTok{Scope} \DataTypeTok{Locked}\NormalTok{) }\DataTypeTok{TInt}
\NormalTok{peakOffsetLocked }\OtherTok{=} \DataTypeTok{EVar}\NormalTok{ (}\DataTypeTok{IS}\NormalTok{ (}\DataTypeTok{IS} \DataTypeTok{IZ}\NormalTok{))}
\end{Highlighting}
\end{Shaded}

And the \texttt{Locked}-specific locals:

\begin{Shaded}
\begin{Highlighting}[]
\CommentTok{{-}{-} source: https://github.com/mstksg/inCode/tree/master/code{-}samples/typed{-}sm{-}lc/MachineStage5.hs\#L173{-}L180}

\OtherTok{lockFreq ::} \DataTypeTok{Expr}\NormalTok{ (}\DataTypeTok{Scope} \DataTypeTok{Locked}\NormalTok{) }\DataTypeTok{TInt}
\NormalTok{lockFreq }\OtherTok{=} \DataTypeTok{EVar}\NormalTok{ (}\DataTypeTok{IS}\NormalTok{ (}\DataTypeTok{IS}\NormalTok{ (}\DataTypeTok{IS} \DataTypeTok{IZ}\NormalTok{)))}

\OtherTok{lockRssi ::} \DataTypeTok{Expr}\NormalTok{ (}\DataTypeTok{Scope} \DataTypeTok{Locked}\NormalTok{) }\DataTypeTok{TInt}
\NormalTok{lockRssi }\OtherTok{=} \DataTypeTok{EVar}\NormalTok{ (}\DataTypeTok{IS}\NormalTok{ (}\DataTypeTok{IS}\NormalTok{ (}\DataTypeTok{IS}\NormalTok{ (}\DataTypeTok{IS} \DataTypeTok{IZ}\NormalTok{))))}

\OtherTok{lockDwell ::} \DataTypeTok{Expr}\NormalTok{ (}\DataTypeTok{Scope} \DataTypeTok{Locked}\NormalTok{) }\DataTypeTok{TInt}
\NormalTok{lockDwell }\OtherTok{=} \DataTypeTok{EVar}\NormalTok{ (}\DataTypeTok{IS}\NormalTok{ (}\DataTypeTok{IS}\NormalTok{ (}\DataTypeTok{IS}\NormalTok{ (}\DataTypeTok{IS}\NormalTok{ (}\DataTypeTok{IS} \DataTypeTok{IZ}\NormalTok{)))))}
\end{Highlighting}
\end{Shaded}

Note that \texttt{lockFreq} is \texttt{EVar\ (IS\ (IS\ (IS\ IZ)))} and
\texttt{rssiLocked} is \texttt{EVar\ IZ}. The types prevent you from
accidentally using \texttt{rssiScanning} (which is
\texttt{Expr\ (Scope\ Scanning)\ TInt}) inside a \texttt{Locked} guard. Even
though both reference position 0, they live in different scopes and GHC treats
them as incompatible.

\subsection{A Destination Has to Build the Right
Payload}\label{a-destination-has-to-build-the-right-payload}

When a transition fires, the machine enters a new state. That new state has its
own \texttt{NodeVars}, which means we need to \emph{construct} the correct
payload from what is available in the source state.

Why can't we just use an expression for this? Because the destination payload is
a \emph{record} of multiple fields, each with its own type. We need to produce
\texttt{NodeVars\ Locked\ =\ {[}"freq"\ :::\ TInt,\ "lockRssi"\ :::\ TInt,\ "dwell"\ :::\ TInt{]}},
which is three separate typed values. A single \texttt{Expr} can only produce
one value. What we need is a collection of expressions, one per field, where
each expression's output type matches the corresponding field's type, and each
expression's scope matches the source state.

This is \texttt{Rec} again. In Part 1, we used \texttt{Rec\ EValueField\ vs} to
store evaluated values for each variable in scope. Here we use
\texttt{Rec\ (BuildField\ scope)\ (NodeVars\ to)} to store \emph{expressions}
for each variable in the destination state. Same container, different payload:

\begin{Shaded}
\begin{Highlighting}[]
\CommentTok{{-}{-} source: https://github.com/mstksg/inCode/tree/master/code{-}samples/typed{-}sm{-}lc/MachineStage5.hs\#L124{-}L125}

\KeywordTok{data} \DataTypeTok{BuildField}\OtherTok{ scope ::}\NormalTok{ (}\DataTypeTok{Symbol}\NormalTok{, }\DataTypeTok{Ty}\NormalTok{) }\OtherTok{{-}\textgreater{}} \DataTypeTok{Type} \KeywordTok{where}
  \DataTypeTok{BuildField}\OtherTok{ ::} \DataTypeTok{String} \OtherTok{{-}\textgreater{}} \DataTypeTok{Expr}\NormalTok{ scope t }\OtherTok{{-}\textgreater{}} \DataTypeTok{BuildField}\NormalTok{ scope (name }\OperatorTok{:::}\NormalTok{ t)}

\KeywordTok{data} \DataTypeTok{BuildField}\OtherTok{ scope ::}\NormalTok{ (}\DataTypeTok{Symbol}\NormalTok{, }\DataTypeTok{Ty}\NormalTok{) }\OtherTok{{-}\textgreater{}} \DataTypeTok{Type} \KeywordTok{where}
  \DataTypeTok{BuildField}\OtherTok{ ::} \DataTypeTok{String} \OtherTok{{-}\textgreater{}} \DataTypeTok{Expr}\NormalTok{ scope t }\OtherTok{{-}\textgreater{}} \DataTypeTok{BuildField}\NormalTok{ scope (name }\OperatorTok{:::}\NormalTok{ t)}
\end{Highlighting}
\end{Shaded}

A \texttt{BuildField\ scope} is indexed by a \texttt{(Symbol,\ Ty)} pair (just
like \texttt{EValueField} was). But instead of holding an evaluated
\texttt{EValue\ t}, it holds an \texttt{Expr\ scope\ t}: a typed expression in
the source state's scope that will \emph{produce} a value of the correct type.
The \texttt{Build} wrapper bundles a name (for rendering in diagrams) with the
full \texttt{Rec} of field expressions.

The \texttt{Rec} GADT enforces that we provide exactly one \texttt{BuildField}
per entry in \texttt{NodeVars\ to}, in order, with matching types. If
\texttt{NodeVars\ Locked} has three fields, the \texttt{Rec} must have exactly
three entries. If the first field is \texttt{"freq"\ :::\ TInt}, the first
\texttt{BuildField} must hold an \texttt{Expr\ scope\ TInt}.

For example, transitioning from \texttt{Scanning} to \texttt{Locked} means
building the payload for \texttt{Locked}, which requires
\texttt{{[}"freq"\ :::\ TInt,\ "lockRssi"\ :::\ TInt,\ "dwell"\ :::\ TInt{]}}:

\begin{Shaded}
\begin{Highlighting}[]
\CommentTok{{-}{-} source: https://github.com/mstksg/inCode/tree/master/code{-}samples/typed{-}sm{-}lc/MachineStage5.hs\#L230{-}L238}

\OtherTok{lockOnPeak ::} \DataTypeTok{Build}\NormalTok{ (}\DataTypeTok{Scope} \DataTypeTok{Scanning}\NormalTok{) }\DataTypeTok{Locked}
\NormalTok{lockOnPeak }\OtherTok{=}
  \DataTypeTok{Build}
    \StringTok{"LockOnPeak"}
\NormalTok{    ( }\DataTypeTok{BuildField} \StringTok{"freq"}\NormalTok{ scanNextFreq}
        \OperatorTok{:\&} \DataTypeTok{BuildField} \StringTok{"lockRssi"}\NormalTok{ rssiScanning}
        \OperatorTok{:\&} \DataTypeTok{BuildField} \StringTok{"dwell"}\NormalTok{ scanNextDwell}
        \OperatorTok{:\&} \DataTypeTok{RNil}
\NormalTok{    )}
\end{Highlighting}
\end{Shaded}

The \texttt{BuildField\ "freq"\ scanNextFreq} provides the
\texttt{"freq"\ :::\ TInt} slot using an expression
\texttt{scanNextFreq\ ::\ Expr\ (Scope\ Scanning)\ TInt}. The
\texttt{BuildField\ "lockRssi"\ rssiScanning} grabs the current RSSI from the
scanning scope. And \texttt{BuildField\ "dwell"\ scanNextDwell} increments the
dwell counter.

The helper expressions used in builders compose the same way as guards:

\begin{Shaded}
\begin{Highlighting}[]
\CommentTok{{-}{-} source: https://github.com/mstksg/inCode/tree/master/code{-}samples/typed{-}sm{-}lc/MachineStage5.hs\#L209{-}L216}

\OtherTok{scanNextFreq ::} \DataTypeTok{Expr}\NormalTok{ (}\DataTypeTok{Scope} \DataTypeTok{Scanning}\NormalTok{) }\DataTypeTok{TInt}
\NormalTok{scanNextFreq }\OtherTok{=} \DataTypeTok{EOp} \DataTypeTok{OPlus}\NormalTok{ scanFreq peakOffsetScanning}

\OtherTok{scanNextDwell ::} \DataTypeTok{Expr}\NormalTok{ (}\DataTypeTok{Scope} \DataTypeTok{Scanning}\NormalTok{) }\DataTypeTok{TInt}
\NormalTok{scanNextDwell }\OtherTok{=} \DataTypeTok{EOp} \DataTypeTok{OPlus}\NormalTok{ scanDwell one}
\end{Highlighting}
\end{Shaded}

\texttt{scanNextFreq} adds \texttt{peakOffset} to the current \texttt{freq}.
\texttt{scanNextDwell} increments \texttt{dwell} by one. Both are
\texttt{Expr\ (Scope\ Scanning)\ TInt}, so they can only be used inside builders
or guards where the source state is \texttt{Scanning}.

If you forget a field, GHC will tell you the \texttt{Rec} doesn't match
\texttt{NodeVars\ Locked}. If you add an extra field, same error. If you provide
an expression of the wrong type for one of the fields, the index won't unify.
The shape of the destination payload is enforced entirely by the type.

Here is a simpler builder that constructs \texttt{NodeVars\ Scanning} from
\texttt{Scope\ Idle}:

\begin{Shaded}
\begin{Highlighting}[]
\CommentTok{{-}{-} source: https://github.com/mstksg/inCode/tree/master/code{-}samples/typed{-}sm{-}lc/MachineStage5.hs\#L221{-}L228}

\OtherTok{startScan ::} \DataTypeTok{Build}\NormalTok{ (}\DataTypeTok{Scope} \DataTypeTok{Idle}\NormalTok{) }\DataTypeTok{Scanning}
\NormalTok{startScan }\OtherTok{=}
  \DataTypeTok{Build}
    \StringTok{"StartScan"}
\NormalTok{    ( }\DataTypeTok{BuildField} \StringTok{"freq"}\NormalTok{ baseFreq}
        \OperatorTok{:\&} \DataTypeTok{BuildField} \StringTok{"dwell"}\NormalTok{ zero}
        \OperatorTok{:\&} \DataTypeTok{RNil}
\NormalTok{    )}
\end{Highlighting}
\end{Shaded}

\texttt{Scope\ Idle} is just \texttt{Env} (since
\texttt{NodeVars\ Idle\ =\ {[}{]}}), so the builder can reference \texttt{rssi},
\texttt{noise}, and \texttt{peakOffset} if it wants to, but here it just uses
constants: \texttt{baseFreq\ =\ 2400} and \texttt{zero\ =\ 0}. And since both
are polymorphic in scope (\texttt{Expr\ scope\ TInt}), they work fine in
\texttt{Scope\ Idle}.

And one that goes the other direction, from \texttt{Locked} back to
\texttt{Scanning}:

\begin{Shaded}
\begin{Highlighting}[]
\CommentTok{{-}{-} source: https://github.com/mstksg/inCode/tree/master/code{-}samples/typed{-}sm{-}lc/MachineStage5.hs\#L240{-}L247}

\OtherTok{keepScanning ::} \DataTypeTok{Build}\NormalTok{ (}\DataTypeTok{Scope} \DataTypeTok{Locked}\NormalTok{) }\DataTypeTok{Scanning}
\NormalTok{keepScanning }\OtherTok{=}
  \DataTypeTok{Build}
    \StringTok{"TrackDrift"}
\NormalTok{    ( }\DataTypeTok{BuildField} \StringTok{"freq"}\NormalTok{ lockNextFreq}
        \OperatorTok{:\&} \DataTypeTok{BuildField} \StringTok{"dwell"}\NormalTok{ lockDwell}
        \OperatorTok{:\&} \DataTypeTok{RNil}
\NormalTok{    )}
\end{Highlighting}
\end{Shaded}

This builder constructs
\texttt{NodeVars\ Scanning\ =\ {[}"freq"\ :::\ TInt,\ "dwell"\ :::\ TInt{]}}
using expressions from \texttt{Scope\ Locked}. \texttt{lockNextFreq} adds the
peak offset to the current locked frequency, and \texttt{lockDwell} carries the
existing dwell count forward.

\subsection{One Transition}\label{one-transition}

Now we can put it all together. A transition packages a source state, a trigger,
a destination state, a guard expression, and a destination builder:

\begin{Shaded}
\begin{Highlighting}[]
\CommentTok{{-}{-} source: https://github.com/mstksg/inCode/tree/master/code{-}samples/typed{-}sm{-}lc/MachineStage5.hs\#L132{-}L139}

\KeywordTok{data} \DataTypeTok{Transition}\NormalTok{ env from trigger to }\KeywordTok{where}
  \DataTypeTok{Transition} \OtherTok{::}
    \DataTypeTok{SState}\NormalTok{ from }\OtherTok{{-}\textgreater{}}
    \DataTypeTok{STrigger}\NormalTok{ trigger }\OtherTok{{-}\textgreater{}}
    \DataTypeTok{SState}\NormalTok{ to }\OtherTok{{-}\textgreater{}}
    \DataTypeTok{Expr}\NormalTok{ (env }\OperatorTok{++} \DataTypeTok{NodeVars}\NormalTok{ from) }\DataTypeTok{TBool} \OtherTok{{-}\textgreater{}}
    \DataTypeTok{Build}\NormalTok{ (env }\OperatorTok{++} \DataTypeTok{NodeVars}\NormalTok{ from) to }\OtherTok{{-}\textgreater{}}
    \DataTypeTok{Transition}\NormalTok{ env from trigger to}
\end{Highlighting}
\end{Shaded}

The type indices enforce:

\begin{itemize}
\tightlist
\item
  The guard expression lives in \texttt{Expr\ (env\ ++\ NodeVars\ from)\ TBool}:
  it can only reference variables available in the source state, and must return
  a boolean.
\item
  The builder lives in \texttt{Build\ (env\ ++\ NodeVars\ from)\ to}: it builds
  the destination payload from the source scope.
\end{itemize}

The \texttt{env} parameter is shared across all transitions in a machine. It
represents the external environment (sensor readings, in our case). The type
parameter is kept abstract so that the same \texttt{Transition} type could be
instantiated against different environments for different machine families.

Here is a simple one: from \texttt{Idle}, on \texttt{Start}, transition to
\texttt{Scanning} unconditionally, building the initial scanning payload:

\begin{Shaded}
\begin{Highlighting}[]
\CommentTok{{-}{-} source: https://github.com/mstksg/inCode/tree/master/code{-}samples/typed{-}sm{-}lc/MachineStage5.hs\#L221{-}L228}

\OtherTok{startScan ::} \DataTypeTok{Build}\NormalTok{ (}\DataTypeTok{Scope} \DataTypeTok{Idle}\NormalTok{) }\DataTypeTok{Scanning}
\NormalTok{startScan }\OtherTok{=}
  \DataTypeTok{Build}
    \StringTok{"StartScan"}
\NormalTok{    ( }\DataTypeTok{BuildField} \StringTok{"freq"}\NormalTok{ baseFreq}
        \OperatorTok{:\&} \DataTypeTok{BuildField} \StringTok{"dwell"}\NormalTok{ zero}
        \OperatorTok{:\&} \DataTypeTok{RNil}
\NormalTok{    )}
\end{Highlighting}
\end{Shaded}

Wait, that is the builder we already defined. The actual \emph{transition} that
uses it looks like:

\begin{Shaded}
\begin{Highlighting}[]
\DataTypeTok{Transition} \DataTypeTok{SIdle} \DataTypeTok{SStart} \DataTypeTok{SScanning}\NormalTok{ true startScan}
\end{Highlighting}
\end{Shaded}

The guard is \texttt{true} (always fire), and the builder provides
\texttt{freq\ =\ 2400} and \texttt{dwell\ =\ 0}. Let us trace through the types.
\texttt{Transition} expects:

\begin{itemize}
\tightlist
\item
  \texttt{SState\ from} -\textgreater{} we give \texttt{SIdle}, so
  \texttt{from\ \textasciitilde{}\ Idle}
\item
  \texttt{STrigger\ trigger} -\textgreater{} \texttt{SStart}, so
  \texttt{trigger\ \textasciitilde{}\ Start}
\item
  \texttt{SState\ to} -\textgreater{} \texttt{SScanning}, so
  \texttt{to\ \textasciitilde{}\ Scanning}
\item
  \texttt{Expr\ (env\ ++\ NodeVars\ from)\ TBool} -\textgreater{}
  \texttt{true\ ::\ Expr\ scope\ TBool}, which unifies with any \texttt{scope},
  so this is fine
\item
  \texttt{Build\ (env\ ++\ NodeVars\ from)\ to} -\textgreater{}
  \texttt{startScan\ ::\ Build\ (Scope\ Idle)\ \ \ Scanning}. The type
  \texttt{Scope\ Idle\ =\ Env\ ++\ NodeVars\ Idle\ =\ Env\ ++\ {[}{]}\ =\ Env},
  and \texttt{env\ ++\ NodeVars\ Idle\ =\ env\ ++\ {[}{]}}. This unifies when
  \texttt{env\ \textasciitilde{}\ Env}.
\end{itemize}

Everything checks out. Now, what happens if we try to transition to
\texttt{Locked} using \texttt{startScan}?

\begin{Shaded}
\begin{Highlighting}[]
\NormalTok{ghci}\OperatorTok{\textgreater{}} \OperatorTok{:}\NormalTok{t }\DataTypeTok{Transition} \DataTypeTok{SIdle} \DataTypeTok{SStart} \DataTypeTok{SLocked}\NormalTok{ true startScan}
\OperatorTok{\textless{}}\NormalTok{interactive}\OperatorTok{\textgreater{}} \FunctionTok{error}\OperatorTok{:}
    \DataTypeTok{Couldn\textquotesingle{}t}\NormalTok{ match }\KeywordTok{type} \DataTypeTok{\textquotesingle{}Scanning\textquotesingle{}}\NormalTok{ with }\DataTypeTok{\textquotesingle{}Locked\textquotesingle{}}
      \DataTypeTok{Expected}\OperatorTok{:} \DataTypeTok{Build}\NormalTok{ (env }\OperatorTok{++} \DataTypeTok{NodeVars} \DataTypeTok{Idle}\NormalTok{) }\DataTypeTok{Locked}
        \DataTypeTok{Actual}\OperatorTok{:} \DataTypeTok{Build}\NormalTok{ (}\DataTypeTok{Scope} \DataTypeTok{Idle}\NormalTok{) }\DataTypeTok{Scanning}
\end{Highlighting}
\end{Shaded}

\texttt{startScan\ ::\ Build\ (Scope\ Idle)\ Scanning}, but we need a
\texttt{Build\ (Scope\ Idle)\ Locked}. GHC rejects it because \texttt{Scanning}
is not \texttt{Locked}.

And what if we try to use \texttt{strongEnoughScanning} (which reads
\texttt{freq} and \texttt{dwell}) as the guard for an \texttt{Idle} transition?

\begin{Shaded}
\begin{Highlighting}[]
\NormalTok{ghci}\OperatorTok{\textgreater{}} \OperatorTok{:}\NormalTok{t }\DataTypeTok{Transition} \DataTypeTok{SIdle} \DataTypeTok{SStart} \DataTypeTok{SScanning}\NormalTok{ strongEnoughScanning startScan}
\OperatorTok{\textless{}}\NormalTok{interactive}\OperatorTok{\textgreater{}} \FunctionTok{error}\OperatorTok{:}
    \DataTypeTok{Couldn\textquotesingle{}t}\NormalTok{ match }\KeywordTok{type}\OperatorTok{:}\NormalTok{ \textquotesingle{}[]}
\NormalTok{                   with}\OperatorTok{:}\NormalTok{ [}\StringTok{"freq"} \OperatorTok{:::} \DataTypeTok{TInt}\NormalTok{, }\StringTok{"dwell"} \OperatorTok{:::} \DataTypeTok{TInt}\NormalTok{]}
      \DataTypeTok{Expected}\OperatorTok{:} \DataTypeTok{Build}\NormalTok{ (env }\OperatorTok{++} \DataTypeTok{NodeVars} \DataTypeTok{Idle}\NormalTok{) }\DataTypeTok{Scanning}
        \DataTypeTok{Actual}\OperatorTok{:} \DataTypeTok{Build}\NormalTok{ (}\DataTypeTok{Scope} \DataTypeTok{Idle}\NormalTok{) }\DataTypeTok{Scanning}
\end{Highlighting}
\end{Shaded}

The error is on \texttt{startScan} here, because GHC tries to unify
\texttt{env\ ++\ NodeVars\ Idle} (from the \texttt{Transition} signature) with
\texttt{Scope\ Scanning} (from \texttt{strongEnoughScanning}), which forces
\texttt{env\ \textasciitilde{}\ Env\ ++\ NodeVars\ Scanning}. But then
\texttt{startScan} expects
\texttt{env\ ++\ NodeVars\ Idle\ =\ (Env\ ++\ NodeVars\ Scanning)\ ++\ {[}{]}},
which doesn't reduce to \texttt{Env\ ++\ {[}{]}\ =\ Env}. The lists don't match
up, and GHC reports the residual mismatch.

The takeaway: you cannot mix expressions from different scopes in the same
transition. The type indices propagate from the singletons through the guard and
builder, and any inconsistency shows up as a unification failure.

\subsection{Many Transitions Need a
Wrapper}\label{many-transitions-need-a-wrapper}

We have a problem.
\texttt{Transition\ SIdle\ SStart\ SScanning\ true\ startScan} has type
\texttt{Transition\ Env\ Idle\ Start\ Scanning}. And the peak-lock transition
has type \texttt{Transition\ Env\ Scanning\ Peak\ Locked}. These are different
types. You cannot put them in the same list.

This is the same problem we faced with \texttt{EValue} in Part 1: an
\texttt{EValue\ TInt} and an \texttt{EValue\ TBool} are different types, so we
needed \texttt{SomeValue} to erase the index and store them together. Here we
need to erase \texttt{from}, \texttt{trigger}, and \texttt{to}:

\begin{Shaded}
\begin{Highlighting}[]
\CommentTok{{-}{-} source: https://github.com/mstksg/inCode/tree/master/code{-}samples/typed{-}sm{-}lc/MachineStage5.hs\#L141{-}L144}

\KeywordTok{data} \DataTypeTok{SomeTransition}\NormalTok{ env }\KeywordTok{where}
  \DataTypeTok{SomeTransition}\OtherTok{ ::} \DataTypeTok{Transition}\NormalTok{ env from trigger to }\OtherTok{{-}\textgreater{}} \DataTypeTok{SomeTransition}\NormalTok{ env}

\KeywordTok{newtype} \DataTypeTok{Machine}\NormalTok{ env }\OtherTok{=} \DataTypeTok{Machine}\NormalTok{ [}\DataTypeTok{SomeTransition}\NormalTok{ env]}
\end{Highlighting}
\end{Shaded}

\texttt{SomeTransition\ env} says: ``I am a transition in the given environment,
but I am not telling you which states or trigger I connect. Trust that my
internals are consistent.'' The \texttt{Machine} is then just a list of these
wrapped transitions.

We lose the ability to statically distinguish transitions from each other in the
list, but we \emph{keep} the guarantee that each individual transition is
internally consistent. The existential seals the proof inside: by the time you
pattern-match on
\texttt{SomeTransition\ (Transition\ from\ trigger\ to\ guard\ build)}, you
recover fresh type variables for \texttt{from}, \texttt{trigger}, and
\texttt{to}, and you know that \texttt{guard} and \texttt{build} are consistent
with them, even though you do not know what those types are from the outside.
This is enough for rendering, because the singletons (\texttt{SState\ from},
\texttt{STrigger\ trigger}, \texttt{SState\ to}) are stored inside and give you
back the runtime names when you need them.

\subsection{A Small Machine}\label{a-small-machine}

With all the pieces in hand, let's define a minimal scan controller with five
transitions:

\begin{Shaded}
\begin{Highlighting}[]
\CommentTok{{-}{-} source: https://github.com/mstksg/inCode/tree/master/code{-}samples/typed{-}sm{-}lc/MachineStage5.hs\#L264{-}L302}

\OtherTok{quickScanMachine ::} \DataTypeTok{Machine} \DataTypeTok{Env}
\NormalTok{quickScanMachine }\OtherTok{=}
  \DataTypeTok{Machine}
\NormalTok{    [ }\DataTypeTok{SomeTransition} \OperatorTok{$}
        \DataTypeTok{Transition}
          \DataTypeTok{SIdle}
          \DataTypeTok{SStart}
          \DataTypeTok{SScanning}
\NormalTok{          true}
\NormalTok{          startScan,}
      \DataTypeTok{SomeTransition} \OperatorTok{$}
        \DataTypeTok{Transition}
          \DataTypeTok{SScanning}
          \DataTypeTok{SPeak}
          \DataTypeTok{SLocked}
\NormalTok{          strongEnoughScanning}
\NormalTok{          lockOnPeak,}
      \DataTypeTok{SomeTransition} \OperatorTok{$}
        \DataTypeTok{Transition}
          \DataTypeTok{SScanning}
          \DataTypeTok{STimeout}
          \DataTypeTok{SCoolingDown}
\NormalTok{          scanTooLong}
\NormalTok{          coolFromScanning,}
      \DataTypeTok{SomeTransition} \OperatorTok{$}
        \DataTypeTok{Transition}
          \DataTypeTok{SLocked}
          \DataTypeTok{SDrift}
          \DataTypeTok{SScanning}
\NormalTok{          true}
\NormalTok{          keepScanning,}
      \DataTypeTok{SomeTransition} \OperatorTok{$}
        \DataTypeTok{Transition}
          \DataTypeTok{SCoolingDown}
          \DataTypeTok{SReset}
          \DataTypeTok{SIdle}
\NormalTok{          true}
\NormalTok{          idle}
\NormalTok{    ]}
\end{Highlighting}
\end{Shaded}

Read it as a graph: Idle -Start-\textgreater{} Scanning, Scanning
-Peak/strongEnough-\textgreater{} Locked, Scanning
-Timeout/tooLong-\textgreater{} CoolingDown, Locked -Drift-\textgreater{}
Scanning, CoolingDown -Reset-\textgreater{} Idle. Each transition carries a
typed guard and a typed builder, both verified against the source state's scope.

\subsection{A Less Toy Machine}\label{a-less-toy-machine}

The \texttt{quickScanMachine} is fine for demonstrating the DSL, but it does not
justify the ceremony. Here is a more interesting one with self-transitions,
multiple paths through the same nodes, and timeout behavior at different stages:

\begin{Shaded}
\begin{Highlighting}[]
\CommentTok{{-}{-} source: https://github.com/mstksg/inCode/tree/master/code{-}samples/typed{-}sm{-}lc/MachineStage5.hs\#L304{-}L363}

\OtherTok{dwellScanMachine ::} \DataTypeTok{Machine} \DataTypeTok{Env}
\NormalTok{dwellScanMachine }\OtherTok{=}
  \DataTypeTok{Machine}
\NormalTok{    [ }\DataTypeTok{SomeTransition} \OperatorTok{$}
        \DataTypeTok{Transition}
          \DataTypeTok{SIdle}
          \DataTypeTok{SStart}
          \DataTypeTok{SScanning}
\NormalTok{          true}
\NormalTok{          startScan,}
      \DataTypeTok{SomeTransition} \OperatorTok{$}
        \DataTypeTok{Transition}
          \DataTypeTok{SScanning}
          \DataTypeTok{SPeak}
          \DataTypeTok{SLocked}
\NormalTok{          strongEnoughScanning}
\NormalTok{          lockOnPeak,}
      \DataTypeTok{SomeTransition} \OperatorTok{$}
        \DataTypeTok{Transition}
          \DataTypeTok{SScanning}
          \DataTypeTok{STimeout}
          \DataTypeTok{SCoolingDown}
\NormalTok{          scanTooLong}
\NormalTok{          coolFromScanning,}
      \DataTypeTok{SomeTransition} \OperatorTok{$}
        \DataTypeTok{Transition}
          \DataTypeTok{SLocked}
          \DataTypeTok{SPeak}
          \DataTypeTok{SLocked}
\NormalTok{          strongEnoughLocked}
\NormalTok{          ( }\DataTypeTok{Build}
              \StringTok{"RefreshLock"}
\NormalTok{              ( }\DataTypeTok{BuildField} \StringTok{"freq"}\NormalTok{ lockFreq}
                  \OperatorTok{:\&} \DataTypeTok{BuildField} \StringTok{"lockRssi"}\NormalTok{ rssiLocked}
                  \OperatorTok{:\&} \DataTypeTok{BuildField} \StringTok{"dwell"}\NormalTok{ lockNextDwell}
                  \OperatorTok{:\&} \DataTypeTok{RNil}
\NormalTok{              )}
\NormalTok{          ),}
      \DataTypeTok{SomeTransition} \OperatorTok{$}
        \DataTypeTok{Transition}
          \DataTypeTok{SLocked}
          \DataTypeTok{SDrift}
          \DataTypeTok{SScanning}
\NormalTok{          true}
\NormalTok{          keepScanning,}
      \DataTypeTok{SomeTransition} \OperatorTok{$}
        \DataTypeTok{Transition}
          \DataTypeTok{SLocked}
          \DataTypeTok{STimeout}
          \DataTypeTok{SCoolingDown}
\NormalTok{          lockedTooLong}
\NormalTok{          coolFromLock,}
      \DataTypeTok{SomeTransition} \OperatorTok{$}
        \DataTypeTok{Transition}
          \DataTypeTok{SCoolingDown}
          \DataTypeTok{SReset}
          \DataTypeTok{SIdle}
\NormalTok{          true}
\NormalTok{          idle}
\NormalTok{    ]}
\end{Highlighting}
\end{Shaded}

The interesting additions:
\texttt{Locked\ -Peak/strongEnough-\textgreater{}\ Locked} is a self-transition
that refreshes the lock (updates \texttt{lockRssi} and increments
\texttt{dwell}).
\texttt{Locked\ -Timeout/lockedTooLong-\textgreater{}\ CoolingDown} adds a
timeout path out of the locked state. Both share the same expression language,
the same \texttt{Scope\ Locked}, and the same \texttt{Build} mechanism. The
self-transition is particularly satisfying: a
\texttt{Build\ (Scope\ Locked)\ Locked} constructs \texttt{NodeVars\ Locked}
from \texttt{Scope\ Locked}, which \emph{contains} \texttt{NodeVars\ Locked}. So
a self-transition can read its own payload and produce an updated version.

Notice that \texttt{dwellScanMachine} reuses the same named expressions
(\texttt{strongEnoughLocked}, \texttt{lockNextFreq}, \texttt{lockNextDwell})
without any risk of accidentally mixing up which state they belong to. The types
keep them honest.

\subsection{Things GHC Will Not Let Us
Say}\label{things-ghc-will-not-let-us-say}

To make this concrete, here are the exact classes of error that are now
compile-time rejections, with real GHC output from attempting them in ghci.

\textbf{Using a variable from the wrong state:}

We have \texttt{lockRssi\ ::\ Expr\ (Scope\ Locked)\ TInt}. Let's try using it
as a guard while \texttt{Scanning}:

\begin{Shaded}
\begin{Highlighting}[]
\NormalTok{ghci}\OperatorTok{\textgreater{}} \OperatorTok{:}\NormalTok{t }\DataTypeTok{Transition} \DataTypeTok{SScanning} \DataTypeTok{SPeak} \DataTypeTok{SLocked}\NormalTok{ lockRssi lockOnPeak}
\OperatorTok{\textless{}}\NormalTok{interactive}\OperatorTok{\textgreater{}} \FunctionTok{error}\OperatorTok{:}
    \DataTypeTok{Couldn\textquotesingle{}t}\NormalTok{ match }\KeywordTok{type} \DataTypeTok{\textquotesingle{}TInt\textquotesingle{}}\NormalTok{ with }\DataTypeTok{\textquotesingle{}TBool\textquotesingle{}}
      \DataTypeTok{Expected}\OperatorTok{:} \DataTypeTok{Expr}\NormalTok{ (env }\OperatorTok{++} \DataTypeTok{NodeVars} \DataTypeTok{Scanning}\NormalTok{) }\DataTypeTok{TBool}
        \DataTypeTok{Actual}\OperatorTok{:} \DataTypeTok{Expr}\NormalTok{ (}\DataTypeTok{Scope} \DataTypeTok{Locked}\NormalTok{) }\DataTypeTok{TInt}
\end{Highlighting}
\end{Shaded}

Two problems: \texttt{lockRssi} returns \texttt{TInt} (not \texttt{TBool}), and
it lives in \texttt{Scope\ Locked} (not \texttt{Scope\ Scanning}). GHC reports
whichever mismatch it encounters first.

\textbf{A guard that returns the wrong type:}

We have \texttt{scanFreq\ ::\ Expr\ (Scope\ Scanning)\ TInt}. Even though it
lives in the right scope, it is not boolean:

\begin{Shaded}
\begin{Highlighting}[]
\NormalTok{ghci}\OperatorTok{\textgreater{}} \OperatorTok{:}\NormalTok{t }\DataTypeTok{Transition} \DataTypeTok{SScanning} \DataTypeTok{SPeak} \DataTypeTok{SLocked}\NormalTok{ scanFreq lockOnPeak}
\OperatorTok{\textless{}}\NormalTok{interactive}\OperatorTok{\textgreater{}} \FunctionTok{error}\OperatorTok{:}
    \DataTypeTok{Couldn\textquotesingle{}t}\NormalTok{ match }\KeywordTok{type} \DataTypeTok{\textquotesingle{}TInt\textquotesingle{}}\NormalTok{ with }\DataTypeTok{\textquotesingle{}TBool\textquotesingle{}}
      \DataTypeTok{Expected}\OperatorTok{:} \DataTypeTok{Expr}\NormalTok{ (env }\OperatorTok{++} \DataTypeTok{NodeVars} \DataTypeTok{Scanning}\NormalTok{) }\DataTypeTok{TBool}
        \DataTypeTok{Actual}\OperatorTok{:} \DataTypeTok{Expr}\NormalTok{ (}\DataTypeTok{Scope} \DataTypeTok{Scanning}\NormalTok{) }\DataTypeTok{TInt}
\end{Highlighting}
\end{Shaded}

The guard slot demands \texttt{TBool}. An integer expression does not fit.

\textbf{A builder with the wrong destination:}

We have \texttt{startScan\ ::\ Build\ (Scope\ Idle)\ Scanning}. Let's try using
it to transition into \texttt{Locked}:

\begin{Shaded}
\begin{Highlighting}[]
\NormalTok{ghci}\OperatorTok{\textgreater{}} \OperatorTok{:}\NormalTok{t }\DataTypeTok{Transition} \DataTypeTok{SIdle} \DataTypeTok{SStart} \DataTypeTok{SLocked}\NormalTok{ true startScan}
\OperatorTok{\textless{}}\NormalTok{interactive}\OperatorTok{\textgreater{}} \FunctionTok{error}\OperatorTok{:}
    \DataTypeTok{Couldn\textquotesingle{}t}\NormalTok{ match }\KeywordTok{type} \DataTypeTok{\textquotesingle{}Scanning\textquotesingle{}}\NormalTok{ with }\DataTypeTok{\textquotesingle{}Locked\textquotesingle{}}
      \DataTypeTok{Expected}\OperatorTok{:} \DataTypeTok{Build}\NormalTok{ (env }\OperatorTok{++} \DataTypeTok{NodeVars} \DataTypeTok{Idle}\NormalTok{) }\DataTypeTok{Locked}
        \DataTypeTok{Actual}\OperatorTok{:} \DataTypeTok{Build}\NormalTok{ (}\DataTypeTok{Scope} \DataTypeTok{Idle}\NormalTok{) }\DataTypeTok{Scanning}
\end{Highlighting}
\end{Shaded}

The builder constructs \texttt{NodeVars\ Scanning}, but the transition's
destination is \texttt{Locked}. \texttt{Scanning} is not \texttt{Locked}.

The important thing here is not that this is ergonomic. It is absolutely not
ergonomic. Writing \texttt{EVar\ (IS\ (IS\ (IS\ IZ)))} to say \texttt{freq} is
painful and error-prone in its own way (though Part 1's \texttt{ListIx} helper
takes the edge off). The important thing is that the illegal versions are not
values. If either machine's \texttt{Peak} guard accidentally returned a
\texttt{TInt}, or if a transition to \texttt{Locked} forgot to construct
\texttt{lockRssi}, there would be no JavaScript to debug because the Haskell
program that generates the JavaScript would not compile.

I want to dwell on this for a moment, because it is genuinely different from
``writing good tests.'' Tests check specific scenarios. Types check \emph{all}
scenarios simultaneously. When I add a new field to \texttt{NodeVars\ Locked},
every builder that targets \texttt{Locked} immediately becomes a type error
until I provide the new field. Every single one. I do not need to remember which
transitions go to \texttt{Locked}. I do not need to grep. GHC will enumerate
them for me in the form of errors. This is the kind of property you cannot get
from testing, no matter how thorough, because tests are intensional (they check
what you wrote) while types are extensional (they check what you \emph{didn't}
write).

\subsection{Two Machines, One Language}\label{two-machines-one-language}

\begin{Shaded}
\begin{Highlighting}[]
\CommentTok{{-}{-} source: https://github.com/mstksg/inCode/tree/master/code{-}samples/typed{-}sm{-}lc/MachineStage5.hs\#L365{-}L369}

\OtherTok{machines ::}\NormalTok{ [(}\DataTypeTok{String}\NormalTok{, }\DataTypeTok{Machine} \DataTypeTok{Env}\NormalTok{)]}
\NormalTok{machines }\OtherTok{=}
\NormalTok{  [ (}\StringTok{"Quick scan"}\NormalTok{, quickScanMachine),}
\NormalTok{    (}\StringTok{"Dwell scan"}\NormalTok{, dwellScanMachine)}
\NormalTok{  ]}
\end{Highlighting}
\end{Shaded}

One generated machine is a demo. Two generated machines show that we built a
\emph{language}. \texttt{quickScanMachine} and \texttt{dwellScanMachine} have
different numbers of transitions, different graph topologies, and different
self-transition behaviors, but they share the same
\texttt{Transition}/\texttt{Build}/\texttt{Expr} types, the same
\texttt{ScanState} vocabulary, and the same rendering backends. The DSL is
general enough to describe an entire family of machines, not just one.

And because the machines are just values (lists of existentially-wrapped
transition records), we can store them, pass them around, and process them with
any function that knows how to walk the structure. There is no metaprogramming,
no Template Haskell, no code generation step. The machines are Haskell values
that happen to carry proofs of their own consistency.

\section{Rendering}\label{rendering}

The first backend does not need to run the machine at all. Because the machine
is data (a list of typed transition records), we can walk the structure and
render it as a graph.

\subsection{The Expression Renderer}\label{the-expression-renderer}

Before we can render transitions, we need to render guard expressions as
human-readable labels. This reuses the same \texttt{Rec\ NameField} /
\texttt{indexRec} pattern from Part 1's pretty-printer: we carry a record of
variable names indexed by the scope, and when we hit an \texttt{EVar\ i}, we
\texttt{indexRec\ i\ names} to get the string name at that position.

\begin{Shaded}
\begin{Highlighting}[]
\CommentTok{{-}{-} source: https://github.com/mstksg/inCode/tree/master/code{-}samples/typed{-}sm{-}lc/MachineStage5.hs\#L77{-}L78}

\KeywordTok{data} \DataTypeTok{NameField}\OtherTok{ ::}\NormalTok{ (}\DataTypeTok{Symbol}\NormalTok{, }\DataTypeTok{Ty}\NormalTok{) }\OtherTok{{-}\textgreater{}} \DataTypeTok{Type} \KeywordTok{where}
  \DataTypeTok{NameField}\OtherTok{ ::} \DataTypeTok{String} \OtherTok{{-}\textgreater{}} \DataTypeTok{NameField}\NormalTok{ field}
\end{Highlighting}
\end{Shaded}

A \texttt{NameField} is indexed by a \texttt{(Symbol,\ Ty)} pair but just holds
a \texttt{String}. It is the rendering analog of \texttt{EValueField}: where
\texttt{EValueField} held a runtime value for evaluation, \texttt{NameField}
holds a display name for rendering.

We define name records for each state's scope:

\begin{Shaded}
\begin{Highlighting}[]
\CommentTok{{-}{-} source: https://github.com/mstksg/inCode/tree/master/code{-}samples/typed{-}sm{-}lc/MachineStage5.hs\#L80{-}L115}

\OtherTok{envNames ::} \DataTypeTok{Rec} \DataTypeTok{NameField} \DataTypeTok{Env}
\NormalTok{envNames }\OtherTok{=}
  \DataTypeTok{NameField} \StringTok{"rssi"}
    \OperatorTok{:\&} \DataTypeTok{NameField} \StringTok{"noise"}
    \OperatorTok{:\&} \DataTypeTok{NameField} \StringTok{"peakOffset"}
    \OperatorTok{:\&} \DataTypeTok{RNil}

\OtherTok{scanningNames ::} \DataTypeTok{Rec} \DataTypeTok{NameField}\NormalTok{ (}\DataTypeTok{Scope} \DataTypeTok{Scanning}\NormalTok{)}
\NormalTok{scanningNames }\OtherTok{=}
  \DataTypeTok{NameField} \StringTok{"rssi"}
    \OperatorTok{:\&} \DataTypeTok{NameField} \StringTok{"noise"}
    \OperatorTok{:\&} \DataTypeTok{NameField} \StringTok{"peakOffset"}
    \OperatorTok{:\&} \DataTypeTok{NameField} \StringTok{"freq"}
    \OperatorTok{:\&} \DataTypeTok{NameField} \StringTok{"dwell"}
    \OperatorTok{:\&} \DataTypeTok{RNil}

\OtherTok{lockedNames ::} \DataTypeTok{Rec} \DataTypeTok{NameField}\NormalTok{ (}\DataTypeTok{Scope} \DataTypeTok{Locked}\NormalTok{)}
\NormalTok{lockedNames }\OtherTok{=}
  \DataTypeTok{NameField} \StringTok{"rssi"}
    \OperatorTok{:\&} \DataTypeTok{NameField} \StringTok{"noise"}
    \OperatorTok{:\&} \DataTypeTok{NameField} \StringTok{"peakOffset"}
    \OperatorTok{:\&} \DataTypeTok{NameField} \StringTok{"freq"}
    \OperatorTok{:\&} \DataTypeTok{NameField} \StringTok{"lockRssi"}
    \OperatorTok{:\&} \DataTypeTok{NameField} \StringTok{"dwell"}
    \OperatorTok{:\&} \DataTypeTok{RNil}

\OtherTok{coolingDownNames ::} \DataTypeTok{Rec} \DataTypeTok{NameField}\NormalTok{ (}\DataTypeTok{Scope} \DataTypeTok{CoolingDown}\NormalTok{)}
\NormalTok{coolingDownNames }\OtherTok{=}
  \DataTypeTok{NameField} \StringTok{"rssi"}
    \OperatorTok{:\&} \DataTypeTok{NameField} \StringTok{"noise"}
    \OperatorTok{:\&} \DataTypeTok{NameField} \StringTok{"peakOffset"}
    \OperatorTok{:\&} \DataTypeTok{NameField} \StringTok{"lastFreq"}
    \OperatorTok{:\&} \DataTypeTok{RNil}
\end{Highlighting}
\end{Shaded}

And a dispatch function that selects the right one given a state singleton:

\begin{Shaded}
\begin{Highlighting}[]
\CommentTok{{-}{-} source: https://github.com/mstksg/inCode/tree/master/code{-}samples/typed{-}sm{-}lc/MachineStage5.hs\#L117{-}L122}

\OtherTok{stateNames ::} \DataTypeTok{SState}\NormalTok{ s }\OtherTok{{-}\textgreater{}} \DataTypeTok{Rec} \DataTypeTok{NameField}\NormalTok{ (}\DataTypeTok{Scope}\NormalTok{ s)}
\NormalTok{stateNames }\OtherTok{=}\NormalTok{ \textbackslash{}}\KeywordTok{case}
  \DataTypeTok{SIdle} \OtherTok{{-}\textgreater{}}\NormalTok{ idleNames}
  \DataTypeTok{SScanning} \OtherTok{{-}\textgreater{}}\NormalTok{ scanningNames}
  \DataTypeTok{SLocked} \OtherTok{{-}\textgreater{}}\NormalTok{ lockedNames}
  \DataTypeTok{SCoolingDown} \OtherTok{{-}\textgreater{}}\NormalTok{ coolingDownNames}
\end{Highlighting}
\end{Shaded}

Pattern matching on \texttt{SIdle} refines the return type to
\texttt{Rec\ NameField\ (Scope\ Idle)}, which is exactly what
\texttt{renderExprDot} will expect when rendering a guard from the \texttt{Idle}
state. The refinement is automatic: \texttt{stateNames\ SScanning} returns
\texttt{Rec\ NameField\ (Scope\ Scanning)}.

The expression renderer itself is straightforward:

\begin{Shaded}
\begin{Highlighting}[]
\CommentTok{{-}{-} source: https://github.com/mstksg/inCode/tree/master/code{-}samples/typed{-}sm{-}lc/MachineStage5.hs\#L408{-}L427}

\OtherTok{renderExprDot ::} \DataTypeTok{Rec} \DataTypeTok{NameField}\NormalTok{ vs }\OtherTok{{-}\textgreater{}} \DataTypeTok{Expr}\NormalTok{ vs t }\OtherTok{{-}\textgreater{}} \DataTypeTok{String}
\NormalTok{renderExprDot names }\OtherTok{=}\NormalTok{ \textbackslash{}}\KeywordTok{case}
  \DataTypeTok{EPrim}\NormalTok{ (}\DataTypeTok{PInt}\NormalTok{ n) }\OtherTok{{-}\textgreater{}} \FunctionTok{show}\NormalTok{ n}
  \DataTypeTok{EPrim}\NormalTok{ (}\DataTypeTok{PBool}\NormalTok{ b) }\OtherTok{{-}\textgreater{}} \KeywordTok{if}\NormalTok{ b }\KeywordTok{then} \StringTok{"true"} \KeywordTok{else} \StringTok{"false"}
  \DataTypeTok{EPrim}\NormalTok{ (}\DataTypeTok{PString}\NormalTok{ s) }\OtherTok{{-}\textgreater{}} \FunctionTok{show}\NormalTok{ s}
  \DataTypeTok{EVar}\NormalTok{ i }\OtherTok{{-}\textgreater{}} \KeywordTok{case}\NormalTok{ indexRec i names }\KeywordTok{of}
    \DataTypeTok{NameField}\NormalTok{ n }\OtherTok{{-}\textgreater{}}\NormalTok{ n}
  \DataTypeTok{EOp}\NormalTok{ o x y }\OtherTok{{-}\textgreater{}}
\NormalTok{    parens }\OperatorTok{$}
\NormalTok{      renderExprDot names x}
        \OperatorTok{++} \StringTok{" "}
        \OperatorTok{++}\NormalTok{ renderOp o}
        \OperatorTok{++} \StringTok{" "}
        \OperatorTok{++}\NormalTok{ renderExprDot names y}
  \DataTypeTok{ERecord}\NormalTok{ \{\} }\OtherTok{{-}\textgreater{}} \StringTok{"\textless{}record\textgreater{}"}
  \DataTypeTok{EAccess}\NormalTok{ \{\} }\OtherTok{{-}\textgreater{}} \StringTok{"\textless{}field\textgreater{}"}
  \DataTypeTok{EChoice}\NormalTok{ \{\} }\OtherTok{{-}\textgreater{}} \StringTok{"\textless{}choice\textgreater{}"}
  \DataTypeTok{ECase}\NormalTok{ \{\} }\OtherTok{{-}\textgreater{}} \StringTok{"\textless{}case\textgreater{}"}
  \DataTypeTok{ELambda}\NormalTok{ \{\} }\OtherTok{{-}\textgreater{}} \StringTok{"\textless{}function\textgreater{}"}
  \DataTypeTok{EApply}\NormalTok{ \{\} }\OtherTok{{-}\textgreater{}} \StringTok{"\textless{}apply\textgreater{}"}
\end{Highlighting}
\end{Shaded}

\texttt{EVar\ i} looks up the name at index \texttt{i} in the names record.
\texttt{EOp} renders infix with parentheses. Primitives render as literals. The
remaining cases (records, sums, lambdas) are not used in our machine guards, so
they render as placeholder strings. In a production system you would handle them
properly, but for SDR scan guards we only need arithmetic and comparisons.

\subsection{The Transition Renderer}\label{the-transition-renderer}

\begin{Shaded}
\begin{Highlighting}[]
\CommentTok{{-}{-} source: https://github.com/mstksg/inCode/tree/master/code{-}samples/typed{-}sm{-}lc/MachineStage5.hs\#L371{-}L406}

\OtherTok{compileDot ::}\NormalTok{ [(}\DataTypeTok{String}\NormalTok{, }\DataTypeTok{Machine} \DataTypeTok{Env}\NormalTok{)] }\OtherTok{{-}\textgreater{}} \DataTypeTok{String}
\NormalTok{compileDot xs }\OtherTok{=}
\NormalTok{  block}
\NormalTok{    [ }\StringTok{"digraph TypedSmLc \{"}\NormalTok{,}
      \StringTok{"  rankdir=LR;"}\NormalTok{,}
      \StringTok{"  node [shape=box];"}\NormalTok{,}
\NormalTok{      indent }\DecValTok{2}\NormalTok{ (intercalate }\StringTok{"\textbackslash{}n"}\NormalTok{ (}\FunctionTok{map}\NormalTok{ (}\FunctionTok{uncurry}\NormalTok{ renderMachineDot) xs)),}
      \StringTok{"\}"}
\NormalTok{    ]}

\OtherTok{renderMachineDot ::} \DataTypeTok{String} \OtherTok{{-}\textgreater{}} \DataTypeTok{Machine} \DataTypeTok{Env} \OtherTok{{-}\textgreater{}} \DataTypeTok{String}
\NormalTok{renderMachineDot name (}\DataTypeTok{Machine}\NormalTok{ transitions) }\OtherTok{=}
\NormalTok{  block}
\NormalTok{    [ }\StringTok{"subgraph cluster\_"} \OperatorTok{++}\NormalTok{ dotId name }\OperatorTok{++} \StringTok{" \{"}\NormalTok{,}
      \StringTok{"  label="} \OperatorTok{++}\NormalTok{ jsString name }\OperatorTok{++} \StringTok{";"}\NormalTok{,}
\NormalTok{      indent }\DecValTok{2}\NormalTok{ (intercalate }\StringTok{"\textbackslash{}n"}\NormalTok{ (}\FunctionTok{map}\NormalTok{ (renderSomeTransitionDot name) transitions)),}
      \StringTok{"\}"}
\NormalTok{    ]}

\OtherTok{renderSomeTransitionDot ::} \DataTypeTok{String} \OtherTok{{-}\textgreater{}} \DataTypeTok{SomeTransition} \DataTypeTok{Env} \OtherTok{{-}\textgreater{}} \DataTypeTok{String}
\NormalTok{renderSomeTransitionDot name (}\DataTypeTok{SomeTransition}\NormalTok{ t) }\OtherTok{=}\NormalTok{ renderTransitionDot name t}

\OtherTok{renderTransitionDot ::} \DataTypeTok{String} \OtherTok{{-}\textgreater{}} \DataTypeTok{Transition} \DataTypeTok{Env}\NormalTok{ from trigger to }\OtherTok{{-}\textgreater{}} \DataTypeTok{String}
\NormalTok{renderTransitionDot name (}\DataTypeTok{Transition}\NormalTok{ from trigger to guard build) }\OtherTok{=}
\NormalTok{  dotNode name (stateName from)}
    \OperatorTok{++} \StringTok{" {-}\textgreater{} "}
    \OperatorTok{++}\NormalTok{ dotNode name (stateName to)}
    \OperatorTok{++} \StringTok{" [label="}
    \OperatorTok{++}\NormalTok{ jsString}
\NormalTok{      ( triggerName trigger}
          \OperatorTok{++} \StringTok{" ["}
          \OperatorTok{++}\NormalTok{ renderExprDot (stateNames from) guard}
          \OperatorTok{++} \StringTok{"] / "}
          \OperatorTok{++}\NormalTok{ buildName build}
\NormalTok{      )}
    \OperatorTok{++} \StringTok{"];"}
\end{Highlighting}
\end{Shaded}

The key function is \texttt{renderTransitionDot}: it pattern-matches on a
\texttt{Transition\ env\ from\ trigger\ to}, extracts the singleton state names
via \texttt{stateName\ from} and \texttt{stateName\ to}, and renders the guard
expression using \texttt{renderExprDot\ (stateNames\ from)\ guard}. The
\texttt{stateNames\ from} call uses the \texttt{from} singleton to select the
correct name record, which is type-compatible with the guard because the
\texttt{Transition} type guarantees that the guard lives in
\texttt{env\ ++\ NodeVars\ from}.

The edge labels in the graph come directly from the typed structure. If you
change a guard expression, the diagram updates. If you rename a state, the
labels update. There is no separate ``documentation'' to keep in sync with the
code, because the diagram \emph{is} the code.

\subsection{The Output}\label{the-output}

Running the program with \texttt{-\/-dot} gives us:

\begin{verbatim}
$ runghc MachineStage5.hs --dot
digraph TypedSmLc {
  rankdir=LR;
  node [shape=box];
  subgraph cluster_Quick_scan {
    label="Quick scan";
    Quick_scan_Idle -> Quick_scan_Scanning [label="Start [true] / StartScan"];
    Quick_scan_Scanning -> Quick_scan_Locked [label="Peak [((noise + 4) <= rssi)] / LockOnPeak"];
    Quick_scan_Scanning -> Quick_scan_CoolingDown [label="Timeout [(3 <= dwell)] / CoolScanner"];
    Quick_scan_Locked -> Quick_scan_Scanning [label="Drift [true] / TrackDrift"];
    Quick_scan_CoolingDown -> Quick_scan_Idle [label="Reset [true] / ResetToIdle"];
  }
  subgraph cluster_Dwell_scan {
    label="Dwell scan";
    Dwell_scan_Idle -> Dwell_scan_Scanning [label="Start [true] / StartScan"];
    Dwell_scan_Scanning -> Dwell_scan_Locked [label="Peak [((noise + 4) <= rssi)] / LockOnPeak"];
    Dwell_scan_Scanning -> Dwell_scan_CoolingDown [label="Timeout [(3 <= dwell)] / CoolScanner"];
    Dwell_scan_Locked -> Dwell_scan_Locked [label="Peak [((noise + 4) <= rssi)] / RefreshLock"];
    Dwell_scan_Locked -> Dwell_scan_Scanning [label="Drift [true] / TrackDrift"];
    Dwell_scan_Locked -> Dwell_scan_CoolingDown [label="Timeout [(3 <= dwell)] / CoolLock"];
    Dwell_scan_CoolingDown -> Dwell_scan_Idle [label="Reset [true] / ResetToIdle"];
  }
}
\end{verbatim}

That gives us a rendered Graphviz view of the two machines:
(\href{/img/entries/typed-sm-lc/machines.png}{full size here})

\begin{figure}
\centering
\pandocbounded{\includegraphics[keepaspectratio]{/img/entries/typed-sm-lc/machines.png}}
\caption{Graphviz rendering of the typed SDR scan machines}
\end{figure}

The generated DOT file is available in the code samples:

The \texttt{Peak} trigger appears in both machines but with different guard
conditions: \texttt{quickScanMachine} only has \texttt{Peak} from
\texttt{Scanning}, while \texttt{dwellScanMachine} also has a \texttt{Peak}
self-transition on \texttt{Locked} (the lock refresh). The guard labels
(\texttt{(noise\ +\ 4)\ \textless{}=\ rssi}) are generated directly by the
expression renderer, not hand-written strings. If we changed \texttt{threshold}
from \texttt{4} to \texttt{6}, the diagram would automatically read
\texttt{(noise\ +\ 6)\ \textless{}=\ rssi} without touching any rendering code.

\subsection{What Else Can We Do With This?}\label{what-else-can-we-do-with-this}

The same \texttt{Machine\ Env} value that produced a Graphviz diagram could also
drive:

\begin{itemize}
\tightlist
\item
  A Haskell simulation: pattern-match on transitions, \texttt{eval} the guards
  against concrete sensor values, and step the machine forward. Because
  \texttt{eval} from Part 1 is total, the simulation cannot crash on a
  well-typed guard.
\item
  Test-case generation: enumerate reachable states by trying all triggers at
  each state and evaluating guards symbolically.
\item
  Documentation tables: render each transition as a row with source, trigger,
  guard (pretty-printed), destination, and payload fields.
\item
  Other language backends: compile the guards and builders to JavaScript, C,
  PureScript, or anything else that can evaluate arithmetic and comparisons.
\end{itemize}

We will do that last one in Part 3: compile both machines to JavaScript and
embed the result as interactive state-machine simulators directly in the page.

\section{The Next Step}\label{the-next-step}

We set out to put our typed expression language inside a larger typed structure,
and we got a state-machine DSL where:

\begin{itemize}
\tightlist
\item
  States are a closed kind, not strings. Misspellings are type errors.
\item
  Each state's payload is declared at the type level via \texttt{NodeVars}.
\item
  Guards are \texttt{Expr\ (Scope\ s)\ TBool}: they can only reference variables
  available in their source state, and they must produce a boolean.
\item
  Builders are \texttt{Rec\ (BuildField\ scope)\ (NodeVars\ to)}: they must
  produce exactly the fields the destination state requires, each with the
  correct type.
\item
  The Graphviz renderer walks the typed structure and produces labels directly
  from the expression pretty-printer, using \texttt{stateNames} to supply
  human-readable variable names at each position.
\end{itemize}

The expression language from Part 1 was useful on its own, but its real power is
as a building block. We threaded it through a larger typed structure and got a
machine description where GHC verifies internal consistency at every joint. The
guard references the right variables. The builder produces the right payload.
The destination exists. All from the same type indices that made \texttt{eval}
total in Part 1.

I think the thing that surprised me most, working with this pattern on real
projects, is how \emph{boring} it becomes once you set it up. Adding a new
transition is mechanical: pick the source and destination singletons, write the
guard in the source scope, write the builder targeting the destination payload,
and GHC either accepts it or tells you exactly what is wrong. There are no
surprises. No ``it compiles but does it actually work?'' anxiety. The machine
either typechecks or it does not, and if it typechecks, you can be confident
that the generated artifacts (diagrams, JavaScript, whatever) will be consistent
with each other.

That boringness is the point. The interesting work is in \emph{designing} the
state space. Once the types are right, filling in the transitions is paint by
numbers. I would rather spend my creative energy on ``what should the states
be'' than on ``did I remember to update the builder for the third transition out
of Locked after adding a field.''

In
\href{/entry/extreme-haskell-typed-state-machines-javascript-compiler.html}{Part
3}, we compile both machines to JavaScript and embed them as live interactive
demos. The same typed expressions that produced the Graphviz labels will produce
the runtime guard and builder functions in the generated code. And because both
backends consume the same \texttt{Machine\ Env} value, they are guaranteed to
agree with each other. One source of truth, multiple outputs. That is the
promise of this whole exercise, and Part 3 is where it lands.

\section{Signoff}\label{signoff}

Hi, thanks for reading! You can reach me via email at
\href{mailto:justin@jle.im}{\nolinkurl{justin@jle.im}}, or at twitter at
\href{https://twitter.com/mstk}{@mstk}! This post and all others are published
under the \href{https://creativecommons.org/licenses/by-nc-nd/3.0/}{CC-BY-NC-ND
3.0} license. Corrections and edits via pull request are welcome and encouraged
at \href{https://github.com/mstksg/inCode}{the source repository}.

If you feel inclined, or this post was particularly helpful for you, why not
consider \href{https://www.patreon.com/justinle/overview}{supporting me on
Patreon}, or a \href{bitcoin:3D7rmAYgbDnp4gp4rf22THsGt74fNucPDU}{BTC donation}?
:)

\textbackslash end\{document\}

\end{document}
